\chapter{สมการชเรอดิงเงอร์และพีชคณิตเชิงเส้นควอนตัม}
\label{chap:schrodinger_equation}
\index{สมการชเรอดิงเงอร์ (Schrödinger Equation)}
\index{สัญกรณ์ดิแรก (Dirac Notation)}
\index{ปริภูมิฮิลเบิร์ต (Hilbert Space)}
\index{ตัวดำเนินการแฮร์มิเชียน (Hermitian Operator)}
\index{ทฤษฎีบทเอห์เรนเฟสต์ (Ehrenfest Theorem)}
\index{กระแสความน่าจะเป็น (Probability Current)}
\index{โครงข่ายประสาทเทียมฟิสิกส์ (PINNs)}

% ==========================================================
% แผนบริหารการสอนประจำบทที่ 3
% ==========================================================
\section*{แผนบริหารการสอนประจำบทที่ 3}
\addcontentsline{toc}{section}{แผนบริหารการสอนประจำบทที่ 3}

\begin{lessonplanbox}
\noindent\textbf{หัวข้อเนื้อหาประจำบท}
\begin{enumerate}
    \item กำเนิดสมการคลื่นควอนตัมและสมการชเรอดิงเงอร์ขึ้นกับเวลา (TDSE)
    \item การแยกตัวแปรและสมการชเรอดิงเงอร์ไม่ขึ้นกับเวลา (TISE)
    \item กฎอนุรักษ์ความน่าจะเป็นและกระแสความน่าจะเป็นใน 1 มิติและ 3 มิติ
    \item โครงสร้างเชิงคณิตศาสตร์ของปริภูมิฮิลเบิร์ตและสัญกรณ์บรา-เคตของดิแรก
    \item ตัวดำเนินการเชิงเส้น ตัวดำเนินการแฮร์มิเชียน และทฤษฎีบทสเปกตรัม
    \item ฐานตั้งฉากปรกติ ความสัมพันธ์ความสมบูรณ์ และการแปลงปริภูมิโมเมนตัม
    \item พีชคณิตคอมมิวเตเตอร์และทฤษฎีบทเอห์เรนเฟสต์ฉบับสมบูรณ์
    \item พรมแดนร่วมสมัย ---- การแก้สมการชเรอดิงเงอร์ด้วยโครงข่ายประสาทเทียมฟิสิกส์ (PINNs)
    \item ห้องปฏิบัติการเสมือนจริง WebAR การจำลองเชิงตัวเลขด้วย Python และการประยุกต์ P-S-C-T
    \item สรุปสาระสำคัญ คำถามเชิงมโนทัศน์ และแบบฝึกหัดท้ายบทฉบับสมบูรณ์
\end{enumerate}

\noindent\textbf{วัตถุประสงค์เชิงพฤติกรรม}
\begin{enumerate}
    \item อธิบายที่มาเชิงฟิสิกส์และอนุพัทธ์สมการชเรอดิงเงอร์ขึ้นกับเวลาและไม่ขึ้นกับเวลาได้อย่างถูกต้องตามหลักการ
    \item พิสูจน์สมการความต่อเนื่องของกระแสความน่าจะเป็นใน 3 มิติ และอธิบายผลของการแปลงเกจแม่เหล็กไฟฟ้าได้
    \item ประยุกต์ใช้สัญกรณ์บรา-เคตของดิแรกในการคำนวณผลคูณภายใน การฉายภาพ และการกระจายสถานะในปริภูมิฮิลเบิร์ตได้
    \item พิสูจน์สมบัติความเป็นจริงของค่าเจาะจงและความตั้งฉากของสถานะเจาะจงของตัวดำเนินการแฮร์มิเชียนได้อย่างรัดกุม
    \item อนุพัทธ์ทฤษฎีบทเอห์เรนเฟสต์ และอธิบายการเชื่อมโยงจากกลศาสตร์ควอนตัมสู่กฎการเคลื่อนที่ข้อที่สองของนิวตันได้
    \item ออกแบบฟังก์ชันการสูญเสีย (Loss Function) ของแบบจำลอง PINNs เพื่อแก้สมการชเรอดิงเงอร์เชิงตัวเลขได้
\end{enumerate}
\end{lessonplanbox}

% ==========================================================
% 3.1 กำเนิดสมการคลื่นควอนตัมและสมการชเรอดิงเงอร์ขึ้นกับเวลา
% ==========================================================
\newpage
\section{กำเนิดสมการคลื่นควอนตัมและสมการชเรอดิงเงอร์ขึ้นกับเวลา}
\index{สมการชเรอดิงเงอร์ขึ้นกับเวลา (TDSE)}
\index{หลักการซ้อนทับควอนตัม (Quantum Superposition Principle)}

ในช่วงปลายคริสต์ศักราช 1925 ถึงต้นปี 1926 แอร์วิน ชเรอดิงเงอร์ (Erwin Schrödinger, 1887--1961) นักฟิสิกส์ทฤษฎีชาวออสเตรีย ได้ต่อยอดสมมติฐานคลื่นสสารของเดอ บรอยล์ โดยมีเป้าหมายเพื่อสร้าง \textbf{สมการคลื่นเชิงอนุพันธ์ (Differential Wave Equation)} ที่สามารถอธิบายพลวัตของคลื่นสสารได้อย่างแม่นยำเฉกเช่นเดียวกับที่สมการของแมกซ์เวลล์อธิบายคลื่นแม่เหล็กไฟฟ้า หรือสมการของนิวตันอธิบายกลศาสตร์ดั้งเดิม

พิจารณาอนุภาคอิสระมวล $m$ เคลื่อนที่ใน 1 มิติด้วยโมเมนตัม $p$ และพลังงาน $E$ ตามแนวคิดของเดอ บรอยล์และพลังค์ อนุภาคนี้สมนัยกับคลื่นระนาบเอกรงค์ (Monochromatic Plane Wave) ซึ่งเขียนแทนด้วยฟังก์ชันเชิงซ้อนได้ดังนี้
\begin{equation}
\Psi(x, t) = A e^{i(kx - \omega t)} = A e^{i(px - Et)/\hbar}
\label{eq:free_plane_wave}
\end{equation}
โดยที่ $k = p/\hbar$ คือเลขคลื่น และ $\omega = E/\hbar$ คือความถี่เชิงมุม

เมื่อหาอนุพันธ์ย่อยเทียบกับเวลา $t$ จะได้
\begin{equation}
\frac{\partial \Psi}{\partial t} = -\frac{iE}{\hbar} \Psi \quad \implies \quad i\hbar \frac{\partial \Psi}{\partial t} = E\Psi
\label{eq:energy_diff}
\end{equation}
และเมื่อหาอนุพันธ์ย่อยเทียบกับตำแหน่ง $x$ สองครั้ง จะได้
\begin{equation}
\frac{\partial \Psi}{\partial x} = \frac{ip}{\hbar} \Psi \quad \implies \quad -i\hbar \frac{\partial \Psi}{\partial x} = p\Psi
\label{eq:momentum_diff}
\end{equation}
\begin{equation}
\frac{\partial^2 \Psi}{\partial x^2} = -\frac{p^2}{\hbar^2} \Psi \quad \implies \quad -\hbar^2 \frac{\partial^2 \Psi}{\partial x^2} = p^2 \Psi
\label{eq:momentum_diff2}
\end{equation}

จากสมการ~\eqref{eq:energy_diff} และ~\eqref{eq:momentum_diff} ชเรอดิงเงอร์ได้นิยามตัวดำเนินการเชิงอนุพันธ์ในกลศาสตร์ควอนตัม ได้แก่ \textbf{ตัวดำเนินการพลังงาน (Energy Operator)} และ \textbf{ตัวดำเนินการโมเมนตัม (Momentum Operator)} ในระบบพิกัดตำแหน่ง ดังนี้
\begin{equation}
\hat{E} = i\hbar \frac{\partial}{\partial t}, \qquad \hat{p}_x = -i\hbar \frac{\partial}{\partial x}
\label{eq:quantum_operators}
\end{equation}

ในกลศาสตร์ดั้งเดิมที่ความเร็วต่ำกว่าแสง พลังงานรวม $E$ ของอนุภาคภายใต้สนามพลังงานศักย์ $V(x, t)$ มีค่าเท่ากับผลรวมของพลังงานจลน์กับพลังงานศักย์
\begin{equation}
E = \frac{p^2}{2m} + V(x, t)
\label{eq:classical_hamiltonian}
\end{equation}

เมื่อแทนปริมาณฟิสิกส์ดั้งเดิมด้วยตัวดำเนินการควอนตัมจากสมการ~\eqref{eq:quantum_operators} แล้วกระทำลงบนฟังก์ชันคลื่น $\Psi(x, t)$ จะได้สมการรากฐานของกลศาสตร์ควอนตัม ซึ่งเรียกว่า \textbf{สมการชเรอดิงเงอร์ขึ้นกับเวลา (Time-Dependent Schrödinger Equation ---- TDSE)} ใน 1 มิติ ดังนี้
\begin{equation}
i\hbar \frac{\partial \Psi(x, t)}{\partial t} = \left[ -\frac{\hbar^2}{2m} \frac{\partial^2}{\partial x^2} + V(x, t) \right] \Psi(x, t)
\label{eq:tdse_1d}
\end{equation}

หากนิยาม \textbf{ตัวดำเนินการฮามิลโทเนียน (Hamiltonian Operator, $\hat{H}$)} ซึ่งแทนพลังงานรวมของระบบ
\begin{equation}
\hat{H} = -\frac{\hbar^2}{2m} \frac{\partial^2}{\partial x^2} + V(x, t) = \frac{\hat{p}_x^2}{2m} + V(x, t)
\label{eq:hamiltonian_def}
\end{equation}
สมการชเรอดิงเงอร์ขึ้นกับเวลาสามารถเขียนในรูปกระชับได้ดังนี้
\begin{equation}
i\hbar \frac{\partial \Psi}{\partial t} = \hat{H}\Psi
\label{eq:tdse_compact}
\end{equation}

สำหรับกรณี 3 มิติ สมการ TDSE ขยายรูปโดยใช้ตัวดำเนินการลาปลาเซียน (Laplacian Operator, $\nabla^2$) ได้เป็น
\begin{equation}
i\hbar \frac{\partial \Psi(\mathbf{r}, t)}{\partial t} = \left[ -\frac{\hbar^2}{2m} \nabla^2 + V(\mathbf{r}, t) \right] \Psi(\mathbf{r}, t)
\label{eq:tdse_3d}
\end{equation}

สมบัติสำคัญที่สุดประการหนึ่งของสมการ TDSE คือ \textbf{ความเป็นสมการเชิงเส้น (Linearity)} นั่นคือ หาก $\Psi_1(x, t)$ และ $\Psi_2(x, t)$ ต่างเป็นผลเฉลยของสมการชเรอดิงเงอร์ การรวมกันเชิงเส้น
\begin{equation}
\Psi(x, t) = c_1 \Psi_1(x, t) + c_2 \Psi_2(x, t)
\label{eq:superposition}
\end{equation}
โดยที่ $c_1, c_2 \in \mathbb{C}$ จะเป็นผลเฉลยของสมการด้วยเสมอ สมบัตินี้เรียกว่า \textbf{หลักการซ้อนทับควอนตัม (Quantum Superposition Principle)} ซึ่งเป็นหัวใจของการประมวลผลข้อมูลควอนตัมและการทำงานของคิวบิต

% ==========================================================
% 3.2 การแยกตัวแปรและสมการชเรอดิงเงอร์ไม่ขึ้นกับเวลา
% ==========================================================
\section{การแยกตัวแปรและสมการชเรอดิงเงอร์ไม่ขึ้นกับเวลา}
\index{สมการชเรอดิงเงอร์ไม่ขึ้นกับเวลา (TISE)}
\index{สถานะคงตัว (Stationary State)}

ในระบบกายภาพจำนวนมาก พลังงานศักย์ไม่เปลี่ยนแปลงตามเวลา กล่าวคือ $V(x, t) = V(x)$ ในสภาวะดังกล่าว เราสามารถแก้สมการชเรอดิงเงอร์ขึ้นกับเวลาได้โดยใช้ \textbf{ระเบียบวิธีแยกตัวแปร (Method of Separation of Variables)} โดยสมมติให้ผลเฉลยอยู่ในรูปผลคูณของฟังก์ชันที่ขึ้นกับตำแหน่งเพียงอย่างเดียว กับฟังก์ชันที่ขึ้นกับเวลาเพียงอย่างเดียว
\begin{equation}
\Psi(x, t) = \psi(x) \phi(t)
\label{eq:sep_vars}
\end{equation}

เมื่อแทนสมการ~\eqref{eq:sep_vars} ลงในสมการ TDSE~\eqref{eq:tdse_1d} จะได้
\begin{equation}
i\hbar \psi(x) \frac{d\phi(t)}{dt} = -\frac{\hbar^2}{2m} \phi(t) \frac{d^2\psi(x)}{dx^2} + V(x) \psi(x) \phi(t)
\end{equation}

หารตลอดด้วย $\psi(x)\phi(t)$ จะได้สมการที่ตัวแปรแยกออกจากกันอย่างสมบูรณ์
\begin{equation}
i\hbar \frac{1}{\phi(t)} \frac{d\phi(t)}{dt} = -\frac{\hbar^2}{2m} \frac{1}{\psi(x)} \frac{d^2\psi(x)}{dx^2} + V(x)
\label{eq:separated_both_sides}
\end{equation}

เนื่องจากทางด้านซ้ายของสมการ~\eqref{eq:separated_both_sides} ขึ้นอยู่กับตัวแปรเวลา $t$ เพียงอย่างเดียว ส่วนทางด้านขวาขึ้นอยู่กับตำแหน่ง $x$ เพียงอย่างเดียว การที่ทั้งสองด้านจะเท่ากันได้ตลอดทุกค่าของ $x$ และ $t$ ทั้งสองด้านต้องมีค่าเท่ากับค่าคงตัวค่าเดียวกัน ซึ่งเรากำหนดให้เป็นค่าคงตัวของการแยกตัวแปร $E$ โดยมีมิติของพลังงาน

\subsection{ส่วนที่ขึ้นกับเวลา}
พิจารณาสมการส่วนเวลา
\begin{equation}
i\hbar \frac{1}{\phi(t)} \frac{d\phi(t)}{dt} = E \quad \implies \quad \frac{d\phi(t)}{dt} = -\frac{iE}{\hbar} \phi(t)
\end{equation}
การอินทิเกรตสมการเชิงอนุพันธ์อันดับหนึ่งนี้ให้ผลเฉลยทันที
\begin{equation}
\phi(t) = e^{-iEt/\hbar} = e^{-i\omega t}
\label{eq:time_part_sol}
\end{equation}
โดยที่ $\omega = E/\hbar$ แสดงว่าวิวัฒนาการตามเวลาของสถานะที่มีพลังงานแน่นอนเป็นเพียงแฟกเตอร์เฟสเชิงซ้อนแบบฮาร์มอนิกเท่านั้น

\subsection{ส่วนที่ขึ้นกับตำแหน่ง ---- สมการชเรอดิงเงอร์ไม่ขึ้นกับเวลา (TISE)}
พิจารณาสมการส่วนตำแหน่ง
\begin{equation}
-\frac{\hbar^2}{2m} \frac{d^2\psi(x)}{dx^2} + V(x)\psi(x) = E\psi(x)
\label{eq:tise_1d}
\end{equation}
สมการนี้เรียกว่า \textbf{สมการชเรอดิงเงอร์ไม่ขึ้นกับเวลา (Time-Independent Schrödinger Equation ---- TISE)} ซึ่งสามารถเขียนในรูปสมการค่าเจาะจงของตัวดำเนินการฮามิลโทเนียนได้ดังนี้
\begin{equation}
\hat{H}\psi(x) = E\psi(x)
\label{eq:tise_operator}
\end{equation}

\subsection{สมบัติของสถานะคงตัว}
ฟังก์ชันคลื่นรวมของสถานะที่มีพลังงานแน่นอน $E$ คือ
\begin{equation}
\Psi(x, t) = \psi(x) e^{-iEt/\hbar}
\label{eq:stationary_full}
\end{equation}
สถานะนี้มีสมบัติพิเศษที่เรียกว่า \textbf{สถานะคงตัว (Stationary State)} ดังนี้
\begin{enumerate}
    \item \textbf{ความหนาแน่นความน่าจะเป็นไม่ขึ้นกับเวลา}
    \begin{equation}
    |\Psi(x, t)|^2 = \Psi^*(x, t)\Psi(x, t) = \left( \psi^*(x) e^{iEt/\hbar} \right) \left( \psi(x) e^{-iEt/\hbar} \right) = |\psi(x)|^2
    \end{equation}
    นั่นคือ โอกาสที่จะพบอนุภาค ณ ตำแหน่งใดๆ ไม่เปลี่ยนแปลงไปตามกาลเวลา
    \item \textbf{ค่าคาดหวังของปริมาณทางกายภาพที่เป็นอิสระจากเวลาไม่เปลี่ยนแปลงตามเวลา}
    \begin{equation}
    \langle Q \rangle = \int_{-\infty}^{\infty} \Psi^*(x, t) \hat{Q} \Psi(x, t) \, dx = \int_{-\infty}^{\infty} \psi^*(x) \hat{Q} \psi(x) \, dx = \text{คงตัว}
    \end{equation}
    \item \textbf{พลังงานรวมมีค่าแน่นอน ปราศจากความไม่แน่นอน} ($\Delta H = 0$)
    \begin{equation}
    \hat{H}^2 \psi = \hat{H}(\hat{H}\psi) = \hat{H}(E\psi) = E(\hat{H}\psi) = E^2\psi \quad \implies \quad \langle H^2 \rangle = E^2, \quad \langle H \rangle = E
    \end{equation}
    ดังนั้น ส่วนเบี่ยงเบนมาตรฐานของพลังงาน $\sigma_E = \sqrt{\langle H^2 \rangle - \langle H \rangle^2} = 0$
\end{enumerate}

ผลเฉลยทั่วไปของสมการ TDSE คือผลรวมเชิงเส้นของสถานะคงตัวทั้งหมดที่มีพลังงาน $E_n$
\begin{equation}
\Psi(x, t) = \sum_{n=1}^{\infty} c_n \psi_n(x) e^{-iE_n t/\hbar}
\label{eq:general_solution_superposition}
\end{equation}
โดยที่สัมประสิทธิ์ $c_n$ ถูกกำหนดจากสภาวะเริ่มต้น $\Psi(x, 0)$ และสอดคล้องกับเงื่อนไขนอร์มอลไลซ์ $\sum_n |c_n|^2 = 1$

% ==========================================================
% 3.3 การอนุรักษ์ความน่าจะเป็นและกระแสความน่าจะเป็น
% ==========================================================
\section{การอนุรักษ์ความน่าจะเป็นและกระแสความน่าจะเป็น}
\index{ความหนาแน่นกระแสความน่าจะเป็น (Probability Current Density)}
\index{สมการความต่อเนื่อง (Continuity Equation)}
\index{การแปลงเกจ (Gauge Transformation)}

การที่ผลรวมของความน่าจะเป็นทั้งหมดในการค้นพบอนุภาคทั่วทั้งเอกภพจะต้องมีค่าเท่ากับ 1 เสมอ ทุกขณะเวลา
\begin{equation}
\int_{-\infty}^{\infty} \rho(x, t) \, dx = \int_{-\infty}^{\infty} |\Psi(x, t)|^2 \, dx = 1
\label{eq:total_prob_conserved}
\end{equation}
นำไปสู่คำถามพื้นฐานว่า กลศาสตร์ควอนตัมรับประกันการอนุรักษ์ความน่าจะเป็นนี้อย่างสอดคล้องในตัวเองหรือไม่ หากเวลาล่วงเลยไป

พิจารณาอัตราการเปลี่ยนแปลงตามเวลาของความหนาแน่นความน่าจะเป็น $\rho(x, t) = \Psi^* \Psi$
\begin{equation}
\frac{\partial \rho}{\partial t} = \frac{\partial}{\partial t}(\Psi^* \Psi) = \frac{\partial \Psi^*}{\partial t}\Psi + \Psi^* \frac{\partial \Psi}{\partial t}
\label{eq:drho_dt_product}
\end{equation}

จากสมการชเรอดิงเงอร์ขึ้นกับเวลา~\eqref{eq:tdse_1d}
\begin{equation}
\frac{\partial \Psi}{\partial t} = \frac{1}{i\hbar} \left( -\frac{\hbar^2}{2m} \frac{\partial^2 \Psi}{\partial x^2} + V \Psi \right) = \frac{i\hbar}{2m} \frac{\partial^2 \Psi}{\partial x^2} - \frac{i}{\hbar} V \Psi
\label{eq:dpsi_dt}
\end{equation}
เมื่อหาสังยุคเชิงซ้อนของสมการ~\eqref{eq:dpsi_dt} (เนื่องจากศักย์ $V(x, t)$ เป็นค่าจริง) จะได้
\begin{equation}
\frac{\partial \Psi^*}{\partial t} = -\frac{i\hbar}{2m} \frac{\partial^2 \Psi^*}{\partial x^2} + \frac{i}{\hbar} V \Psi^*
\label{eq:dpsistar_dt}
\end{equation}

นำสมการ~\eqref{eq:dpsi_dt} และ~\eqref{eq:dpsistar_dt} แทนกลับลงในสมการ~\eqref{eq:drho_dt_product}
\begin{align}
\frac{\partial \rho}{\partial t} &= \left( -\frac{i\hbar}{2m} \frac{\partial^2 \Psi^*}{\partial x^2} + \frac{i}{\hbar} V \Psi^* \right) \Psi + \Psi^* \left( \frac{i\hbar}{2m} \frac{\partial^2 \Psi}{\partial x^2} - \frac{i}{\hbar} V \Psi \right) \nonumber \\
&= \frac{i\hbar}{2m} \left( \Psi^* \frac{\partial^2 \Psi}{\partial x^2} - \frac{\partial^2 \Psi^*}{\partial x^2} \Psi \right) \nonumber \\
&= -\frac{\partial}{\partial x} \left[ \frac{\hbar}{2mi} \left( \Psi^* \frac{\partial \Psi}{\partial x} - \frac{\partial \Psi^*}{\partial x} \Psi \right) \right]
\label{eq:continuity_derivation}
\end{align}

เราจึงนิยาม \textbf{ความหนาแน่นกระแสความน่าจะเป็น (Probability Current Density, $j(x, t)$)} ใน 1 มิติ ดังนี้
\begin{equation}
j(x, t) = \frac{\hbar}{2mi} \left( \Psi^* \frac{\partial \Psi}{\partial x} - \Psi \frac{\partial \Psi^*}{\partial x} \right) = \frac{\hbar}{m} \text{Im}\left( \Psi^* \frac{\partial \Psi}{\partial x} \right)
\label{eq:prob_current_1d}
\end{equation}

จากสมการ~\eqref{eq:continuity_derivation} เราได้ \textbf{สมการความต่อเนื่อง (Continuity Equation)} ซึ่งมีรูปแบบทางคณิตศาสตร์เช่นเดียวกับกฎอนุรักษ์ประจุในทฤษฎีแม่เหล็กไฟฟ้าและการอนุรักษ์มวลในกลศาสตร์ของไหล
\begin{equation}
\frac{\partial \rho}{\partial t} + \frac{\partial j}{\partial x} = 0
\label{eq:continuity_equation}
\end{equation}

ใน 3 มิติ ความหนาแน่นกระแสความน่าจะเป็นเขียนในรูปเวกเตอร์ $\mathbf{j}(\mathbf{r}, t)$ ได้เป็น
\begin{equation}
\mathbf{j}(\mathbf{r}, t) = \frac{\hbar}{2mi} \left( \Psi^* \nabla \Psi - \Psi \nabla \Psi^* \right) = \frac{\hbar}{m} \text{Im}(\Psi^* \nabla \Psi)
\label{eq:prob_current_3d}
\end{equation}
และสมการความต่อเนื่องใน 3 มิติ คือ
\begin{equation}
\frac{\partial \rho}{\partial t} + \nabla \cdot \mathbf{j} = 0
\label{eq:continuity_3d}
\end{equation}

\subsection{กระแสความน่าจะเป็นในสนามแม่เหล็กไฟฟ้าและการแปลงเกจ}

เมื่ออนุภาคมีประจุ $q$ เคลื่อนที่ภายใต้สนามแม่เหล็กไฟฟ้าที่มีเวกเตอร์ศักย์ $\mathbf{A}(\mathbf{r}, t)$ ตัวดำเนินการโมเมนตัมคาโนนิคัล $\hat{\mathbf{p}}$ จะถูกแทนที่ด้วยโมเมนตัมจลนศาสตร์ $\hat{\mathbf{P}} = \hat{\mathbf{p}} - q\mathbf{A}$ ตามหลักการจับคู่แบบมินิมัล (Minimal Coupling) ส่งผลให้ความหนาแน่นกระแสความน่าจะเป็นขยายรูปเป็น

\begin{equation}
\mathbf{j}(\mathbf{r}, t) = \frac{1}{m} \text{Re}\left[ \Psi^* \left( -i\hbar\nabla - q\mathbf{A} \right) \Psi \right] = \frac{\hbar}{m}\text{Im}(\Psi^*\nabla\Psi) - \frac{q}{m}\mathbf{A}|\Psi|^2
\label{eq:prob_current_gauge}
\end{equation}

เมื่อระบบอยู่ภายใต้การแปลงเกจ (Gauge Transformation) $\mathbf{A}' = \mathbf{A} + \nabla \chi$ ฟังก์ชันคลื่นจะเปลี่ยนเฟสตามความสัมพันธ์ $\Psi' = \Psi e^{i q\chi / \hbar}$ ซึ่งทำให้กระแสความน่าจะเป็น $\mathbf{j}' = \mathbf{j}$ มีค่าคงเดิมเสมอ ปรากฏการณ์นี้เรียกว่า \textbf{ความไม่แปรเปลี่ยนตามเกจ (Gauge Invariance)} ซึ่งเป็นรากฐานของทฤษฎีเกจสมัยใหม่และปรากฏการณ์อาฮารอนอฟ-โบห์ม (Aharonov-Bohm Effect)

% ==========================================================
% 3.4 สัญกรณ์ดิแรกและโครงสร้างปริภูมิฮิลเบิร์ต
% ==========================================================
\section{สัญกรณ์ดิแรกและโครงสร้างปริภูมิฮิลเบิร์ต}
\index{สัญกรณ์บรา-เคต (Bra-Ket Notation)}
\index{ปริภูมิฮิลเบิร์ต (Hilbert Space)}
\index{ความสมบูรณ์ของฐาน (Completeness Relation)}

ในคริสต์ศักราช 1930 พอล ดิแรก (Paul Dirac, 1902--1984) นักฟิสิกส์ทฤษฎีชาวอังกฤษ ได้เสนอระบบสัญกรณ์ที่ทรงพลังและงดงามเป็นเลิศในหนังสือ \textit{The Principles of Quantum Mechanics} เรียกว่า \textbf{สัญกรณ์บรา-เคต (Bra-Ket Notation)} ซึ่งช่วยรวมฟิสิกส์กลศาสตร์คลื่นของชเรอดิงเงอร์และกลศาสตร์เมทริกซ์ของไฮเซนเบิร์กเข้าสู่โครงสร้างทางคณิตศาสตร์อันเป็นหนึ่งเดียวบน \textbf{ปริภูมิฮิลเบิร์ต (Hilbert Space, $\mathcal{H}$)}

\subsection{เวกเตอร์เคตและบรา}
ในมุมมองเชิงพีชคณิตเชิงเส้น สถานะควอนตัมของระบบถูกแทนด้วยเวกเตอร์ในปริภูมิเชิงซ้อนมิติอนันต์ เรียกว่า \textbf{เวกเตอร์เคต (Ket Vector)} เขียนแทนด้วยสัญลักษณ์ $\ket{\psi} \in \mathcal{H}$ 

สำหรับทุกๆ เคต $\ket{\psi}$ จะมีเวกเตอร์คู่กันในปริภูมิดูอัล (Dual Space, $\mathcal{H}^*$) เรียกว่า \textbf{เวกเตอร์บรา (Bra Vector)} เขียนแทนด้วยสัญลักษณ์
\begin{equation}
\bra{\psi} = \big( \ket{\psi} \big)^\dagger \in \mathcal{H}^*
\end{equation}
โดยที่เครื่องหมายกริช $(\dagger)$ หมายถึงสังยุคแฮร์มิเชียน (Hermitian Adjoint หรือ Conjugate Transpose)

\subsection{ผลคูณภายในและสมบัติเมทริกซ์}
ผลคูณภายใน (Inner Product) ระหว่างบรา $\bra{\phi}$ และเคต $\ket{\psi}$ ประกบกันเป็นคำว่า Bracket เขียนแทนด้วย $\braket{\phi}{\psi} \in \mathbb{C}$ ซึ่งมีสมบัติพื้นฐานดังนี้
\begin{enumerate}
    \item \textbf{สมบัติสังยุคสมมาตร (Conjugate Symmetry)} ---- $\braket{\phi}{\psi} = \braket{\psi}{\phi}^*$
    \item \textbf{ความเป็นเชิงเส้นในอาร์กิวเมนต์ที่สอง (Linearity in Ket)} ---- $\braket{\phi}{c_1 \psi_1 + c_2 \psi_2} = c_1 \braket{\phi}{\psi_1} + c_2 \braket{\phi}{\psi_2}$
    \item \textbf{ความเป็นแอนติลีเนียร์ในอาร์กิวเมนต์แรก (Antilinearity in Bra)} ---- $\braket{c_1 \phi_1 + c_2 \phi_2}{\psi} = c_1^* \braket{\phi_1}{\psi} + c_2^* \braket{\phi_2}{\psi}$
    \item \textbf{ภาวะบวกแน่นอน (Positive Definiteness)} ---- $\braket{\psi}{\psi} \ge 0$ โดยที่ $\braket{\psi}{\psi} = 0 \iff \ket{\psi} = 0$
\end{enumerate}

\subsection{ความสัมพันธ์ระหว่างฟังก์ชันคลื่นกับสัญกรณ์ดิแรก}
ในปริภูมิฮิลเบิร์ตของตำแหน่ง ให้ $\{\ket{x}\}$ เป็นเซตของสถานะเจาะจงของตำแหน่ง ฟังก์ชันคลื่น $\psi(x)$ ในกลศาสตร์คลื่นแท้จริงแล้วคือสัมประสิทธิ์การฉายภาพของเคตสถานะ $\ket{\psi}$ ลงบนแกนตำแหน่ง $\ket{x}$
\begin{equation}
\psi(x) = \braket{x}{\psi}, \qquad \psi^*(x) = \braket{\psi}{x}
\label{eq:wavefunc_projection}
\end{equation}

ในทำนองเดียวกัน ฟังก์ชันคลื่นในปริภูมิโมเมนตัม $\Phi(p)$ คือการฉายภาพลงบนสถานะโมเมนตัม $\ket{p}$
\begin{equation}
\Phi(p) = \braket{p}{\psi}
\end{equation}

ผลคูณภายในระหว่างสองสถานะ $\ket{\phi}$ และ $\ket{\psi}$ ในการแทนด้วยตำแหน่ง สอดคล้องกับการอินทิเกรตฟังก์ชันคลื่น
\begin{equation}
\braket{\phi}{\psi} = \int_{-\infty}^{\infty} \braket{\phi}{x}\braket{x}{\psi} \, dx = \int_{-\infty}^{\infty} \phi^*(x) \psi(x) \, dx
\end{equation}

\subsection{ความสมบูรณ์และตัวดำเนินการฉายภาพ}
หาก $\{ \ket{u_n} \}$ เป็นเซตฐานตั้งฉากปรกติ (Orthonormal Basis) ของปริภูมิฮิลเบิร์ต ซึ่งสอดคล้องกับเงื่อนไขภาวะตั้งฉากปรกติ $\braket{u_m}{u_n} = \delta_{mn}$ แล้ว สถานะใดๆ $\ket{\psi}$ สามารถกระจายบนฐานนี้ได้เป็น
\begin{equation}
\ket{\psi} = \sum_n c_n \ket{u_n}, \qquad c_n = \braket{u_n}{\psi}
\label{eq:basis_expansion}
\end{equation}

เมื่อนำ $c_n$ แทนกลับเข้าไป
\begin{equation}
\ket{\psi} = \sum_n \ket{u_n}\braket{u_n}{\psi} = \left( \sum_n \ket{u_n}\bra{u_n} \right) \ket{\psi}
\end{equation}
เนื่องจากสมการนี้เป็นจริงสำหรับทุกเคต $\ket{\psi}$ จะได้ \textbf{ความสัมพันธ์ความสมบูรณ์ (Completeness Relation หรือ Resolution of the Identity)}
\begin{equation}
\sum_n \ket{u_n}\bra{u_n} = \hat{I}
\label{eq:completeness_discrete}
\end{equation}
โดยที่ตัวดำเนินการ $\hat{P}_n = \ket{u_n}\bra{u_n}$ คือ \textbf{ตัวดำเนินการฉายภาพ (Projection Operator)} ลงบนสถานะ $\ket{u_n}$

สำหรับตัวแปรต่อเนื่อง เช่น ตำแหน่ง $x$ ความสัมพันธ์ความสมบูรณ์เขียนในรูปการอินทิเกรต
\begin{equation}
\int_{-\infty}^{\infty} \ket{x}\bra{x} \, dx = \hat{I}
\label{eq:completeness_continuous}
\end{equation}

% ==========================================================
% 3.5 ตัวดำเนินการแฮร์มิเชียน ปัญหาค่าเจาะจง และทฤษฎีบทสเปกตรัม
% ==========================================================
\section{ตัวดำเนินการแฮร์มิเชียน ปัญหาค่าเจาะจง และทฤษฎีบทสเปกตรัม}
\index{ตัวดำเนินการแฮร์มิเชียน (Hermitian Operator)}
\index{ค่าเจาะจงและสถานะเจาะจง (Eigenvalue and Eigenstate)}
\index{ทฤษฎีบทสเปกตรัม (Spectral Theorem)}

ในกลศาสตร์ควอนตัม ปริมาณที่สังเกตวัดได้จริงทางกายภาพ (Observables) ทั้งหมด เช่น ตำแหน่ง โมเมนตัม พลังงาน และโมเมนตัมเชิงมุม จะถูกแทนด้วย \textbf{ตัวดำเนินการเชิงเส้นที่เป็นแฮร์มิเชียน (Hermitian Linear Operator)}

\subsection{นิยามและสมบัติของตัวดำเนินการแฮร์มิเชียน}
ตัวดำเนินการ $\hat{A}$ จะเป็นตัวดำเนินการแฮร์มิเชียน ก็ต่อเมื่อมันเท่ากับสังยุคแฮร์มิเชียนของตัวเอง
\begin{equation}
\hat{A}^\dagger = \hat{A}
\label{eq:hermitian_def}
\end{equation}
ในรูปของผลคูณภายใน สำหรับทุกๆ เคต $\ket{\phi}$ และ $\ket{\psi}$ ในปริภูมิฮิลเบิร์ต
\begin{equation}
\braket{\phi}{\hat{A}\psi} = \braket{\hat{A}\phi}{\psi}
\label{eq:hermitian_inner_product}
\end{equation}

\subsection{ทฤษฎีบทที่ 1 ---- ค่าเจาะจงของตัวดำเนินการแฮร์มิเชียนเป็นจำนวนจริงเสมอ}
\textbf{การพิสูจน์} \\
ให้ $\ket{a}$ เป็นสถานะเจาะจงของตัวดำเนินการแฮร์มิเชียน $\hat{A}$ ที่มีค่าเจาะจง $a$
\begin{equation}
\hat{A}\ket{a} = a\ket{a}
\label{eq:eigen_eq1}
\end{equation}
นำบรา $\bra{a}$ มาประกบจากทางซ้าย
\begin{equation}
\bra{a}\hat{A}\ket{a} = a\braket{a}{a}
\label{eq:proof_real_1}
\end{equation}
หาคอนจูเกตทรานสโพสของสมการ~\eqref{eq:eigen_eq1} จะได้
\begin{equation}
\bra{a}\hat{A}^\dagger = a^*\bra{a} \quad \implies \quad \bra{a}\hat{A} = a^*\bra{a}
\end{equation}
นำเคต $\ket{a}$ มาประกบจากทางขวา
\begin{equation}
\bra{a}\hat{A}\ket{a} = a^*\braket{a}{a}
\label{eq:proof_real_2}
\end{equation}
นำสมการ~\eqref{eq:proof_real_1} ลบด้วยสมการ~\eqref{eq:proof_real_2}
\begin{equation}
(a - a^*)\braket{a}{a} = 0
\end{equation}
เนื่องจากเคตสถานะเจาะจงไม่ใช่เวกเตอร์ศูนย์ ดังนั้น $\braket{a}{a} > 0$ จึงบังคับให้
\begin{equation}
a - a^* = 0 \quad \implies \quad a = a^* \in \mathbb{R}
\end{equation}
นั่นคือ ค่าเจาะจง $a$ ต้องเป็นจำนวนจริงอย่างแน่นอน ซึ่งสอดคล้องกับข้อเท็จจริงว่า ผลการวัดทางกายภาพในห้องปฏิบัติการต้องเป็นค่าจริงเสมอ

\subsection{ทฤษฎีบทที่ 2 ---- สถานะเจาะจงที่มีค่าเจาะจงต่างกันจะตั้งฉากกันเสมอ}
\textbf{การพิสูจน์} \\
พิจารณาสถานะเจาะจงสองสถานะ $\ket{a_1}$ และ $\ket{a_2}$ โดยที่ $a_1 \ne a_2$
\begin{align}
\hat{A}\ket{a_1} &= a_1\ket{a_1} \label{eq:ortho_1} \\
\hat{A}\ket{a_2} &= a_2\ket{a_2} \label{eq:ortho_2}
\end{align}
นำ $\bra{a_2}$ ประกบทางซ้ายของสมการ~\eqref{eq:ortho_1}
\begin{equation}
\bra{a_2}\hat{A}\ket{a_1} = a_1\braket{a_2}{a_1}
\label{eq:ortho_lhs}
\end{equation}
จากสมการ~\eqref{eq:ortho_2} เมื่อหาสังยุคแฮร์มิเชียนและใช้ความเป็นจริงของ $a_2$ จะได้ $\bra{a_2}\hat{A} = a_2\bra{a_2}$ นำ $\ket{a_1}$ มาประกบจากทางขวา
\begin{equation}
\bra{a_2}\hat{A}\ket{a_1} = a_2\braket{a_2}{a_1}
\label{eq:ortho_rhs}
\end{equation}
นำสมการ~\eqref{eq:ortho_lhs} ลบด้วยสมการ~\eqref{eq:ortho_rhs}
\begin{equation}
(a_1 - a_2)\braket{a_2}{a_1} = 0
\end{equation}
เนื่องจากโจทย์กำหนดให้ค่าเจาะจงไม่เท่ากัน ($a_1 \ne a_2$ ทำให้ $a_1 - a_2 \ne 0$) จึงสรุปได้ว่า
\begin{equation}
\braket{a_2}{a_1} = 0
\end{equation}
นั่นคือ สถานะเจาะจงทั้งสองตั้งฉากกัน (Orthogonal)

\subsection{ทฤษฎีบทสเปกตรัม (Spectral Theorem)}
จากสมบัติความสมบูรณ์และความเป็นแฮร์มิเชียน ตัวดำเนินการ $\hat{A}$ สามารถเขียนกระจายในรูปของการรวมตัวดำเนินการฉายภาพตาม \textbf{ทฤษฎีบทสเปกตรัม (Spectral Decomposition Theorem)}
\begin{equation}
\hat{A} = \sum_n a_n \ket{a_n}\bra{a_n} = \sum_n a_n \hat{P}_n
\label{eq:spectral_decomposition}
\end{equation}
และฟังก์ชันใดๆ ของตัวดำเนินการ $f(\hat{A})$ สามารถคำนวณได้อย่างง่ายดายผ่าน
\begin{equation}
f(\hat{A}) = \sum_n f(a_n) \ket{a_n}\bra{a_n}
\label{eq:operator_function}
\end{equation}

% ==========================================================
% 3.6 พีชคณิตคอมมิวเตเตอร์และทฤษฎีบทเอห์เรนเฟสต์
% ==========================================================
\section{พีชคณิตคอมมิวเตเตอร์และทฤษฎีบทเอห์เรนเฟสต์}
\index{คอมมิวเตเตอร์ (Commutator)}
\index{ทฤษฎีบทเอห์เรนเฟสต์ (Ehrenfest Theorem)}
\index{หลักความสอดคล้อง (Correspondence Principle)}

\subsection{พีชคณิตคอมมิวเตเตอร์และเอกลักษณ์สำคัญ}
คอมมิวเตเตอร์ระหว่างตัวดำเนินการสองตัว $\hat{A}$ และ $\hat{B}$ นิยามโดย
\begin{equation}
[\hat{A}, \hat{B}] = \hat{A}\hat{B} - \hat{B}\hat{A}
\label{eq:commutator_def}
\end{equation}
เอกลักษณ์พื้นฐานที่สำคัญในพีชคณิตคอมมิวเตเตอร์ ได้แก่
\begin{enumerate}
    \item \textbf{การสลับที่แบบปฏิสมมาตร} ---- $[\hat{A}, \hat{B}] = -[\hat{B}, \hat{A}]$
    \item \textbf{ความเป็นเชิงเส้น} ---- $[\hat{A}, \hat{B} + \hat{C}] = [\hat{A}, \hat{B}] + [\hat{A}, \hat{C}]$
    \item \textbf{กฎอนุพันธ์ผลคูณของไลบ์นิซ}
    \begin{align}
    [\hat{A}\hat{B}, \hat{C}] &= \hat{A}[\hat{B}, \hat{C}] + [\hat{A}, \hat{C}]\hat{B} \\
    [\hat{A}, \hat{B}\hat{C}] &= \hat{B}[\hat{A}, \hat{C}] + [\hat{A}, \hat{B}]\hat{C}
    \end{align}
    \item \textbf{เอกลักษณ์ยาโคบี (Jacobi Identity)}
    \begin{equation}
    [\hat{A}, [\hat{B}, \hat{C}]] + [\hat{B}, [\hat{C}, \hat{A}]] + [\hat{C}, [\hat{A}, \hat{B}]] = 0
    \end{equation}
\end{enumerate}

คอมมิวเตเตอร์คาโนนิคัลระหว่างตำแหน่งและโมเมนตัม $[\hat{x}, \hat{p}_x] = i\hbar \hat{I}$ สามารถขยายไปสู่ฟังก์ชันใดๆ ของตำแหน่ง $f(x)$ และโมเมนตัม $g(p)$ ดังนี้
\begin{equation}
[f(\hat{x}), \hat{p}_x] = i\hbar \frac{\partial f}{\partial x}, \qquad [\hat{x}, g(\hat{p}_x)] = i\hbar \frac{\partial g}{\partial p_x}
\label{eq:functional_commutator}
\end{equation}

\subsection{การอนุพัทธ์ทฤษฎีบทเอห์เรนเฟสต์ฉบับสมบูรณ์}
ทฤษฎีบทเอห์เรนเฟสต์ (Ehrenfest Theorem) เสนอโดย พอล เอห์เรนเฟสต์ ในปี 1927 เป็นสะพานเชื่อมที่งดงามที่สุดระหว่างกลศาสตร์ควอนตัมและกลศาสตร์ดั้งเดิม

พิจารณาอัตราการเปลี่ยนแปลงตามเวลาของค่าคาดหวังของตัวดำเนินการ $\hat{A}$
\begin{equation}
\frac{d}{dt}\langle A \rangle = \frac{d}{dt}\braket{\Psi(t)}{\hat{A}\Psi(t)} = \left(\frac{\partial}{\partial t}\bra{\Psi}\right) \hat{A} \ket{\Psi} + \bra{\Psi} \frac{\partial \hat{A}}{\partial t} \ket{\Psi} + \bra{\Psi} \hat{A} \left(\frac{\partial}{\partial t}\ket{\Psi}\right)
\end{equation}

จากสมการชเรอดิงเงอร์ $i\hbar \frac{\partial}{\partial t}\ket{\Psi} = \hat{H}\ket{\Psi} \implies \frac{\partial}{\partial t}\ket{\Psi} = \frac{1}{i\hbar}\hat{H}\ket{\Psi}$ และสังยุคแฮร์มิเชียน $\frac{\partial}{\partial t}\bra{\Psi} = -\frac{1}{i\hbar}\bra{\Psi}\hat{H}$ แทนค่าลงไป
\begin{align}
\frac{d}{dt}\langle A \rangle &= -\frac{1}{i\hbar}\bra{\Psi}\hat{H}\hat{A}\ket{\Psi} + \left\langle \frac{\partial \hat{A}}{\partial t} \right\rangle + \frac{1}{i\hbar}\bra{\Psi}\hat{A}\hat{H}\ket{\Psi} \nonumber \\
&= \frac{1}{i\hbar}\bra{\Psi} (\hat{A}\hat{H} - \hat{H}\hat{A}) \ket{\Psi} + \left\langle \frac{\partial \hat{A}}{\partial t} \right\rangle
\end{align}

เราจึงได้ \textbf{สมการเอห์เรนเฟสต์ (Ehrenfest Equation)}
\begin{equation}
\frac{d}{dt}\langle A \rangle = \frac{1}{i\hbar} \langle [\hat{A}, \hat{H}] \rangle + \left\langle \frac{\partial \hat{A}}{\partial t} \right\rangle
\label{eq:ehrenfest_master}
\end{equation}

\subsection{การเกิดใหม่ของกฎการเคลื่อนที่ของนิวตัน}
เมื่อประยุกต์สมการเอห์เรนเฟสต์เข้ากับตัวดำเนินการตำแหน่ง $\hat{x}$ ภายใต้ฮามิลโทเนียน $\hat{H} = \frac{\hat{p}^2}{2m} + V(x)$
\begin{equation}
[\hat{x}, \hat{H}] = \left[ \hat{x}, \frac{\hat{p}^2}{2m} \right] = \frac{1}{2m} (\hat{p}[\hat{x}, \hat{p}] + [\hat{x}, \hat{p}]\hat{p}) = \frac{1}{2m} (\hat{p}(i\hbar) + (i\hbar)\hat{p}) = \frac{i\hbar}{m}\hat{p}
\end{equation}
\begin{equation}
\frac{d}{dt}\langle x \rangle = \frac{1}{i\hbar} \left\langle \frac{i\hbar}{m}\hat{p} \right\rangle = \frac{\langle p \rangle}{m}
\label{eq:ehrenfest_x}
\end{equation}

เมื่อประยุกต์เข้ากับตัวดำเนินการโมเมนตัม $\hat{p}$
\begin{equation}
[\hat{p}, \hat{H}] = [\hat{p}, V(x)] = -i\hbar \frac{\partial V}{\partial x}
\end{equation}
\begin{equation}
\frac{d}{dt}\langle p \rangle = \frac{1}{i\hbar} \left\langle -i\hbar \frac{\partial V}{\partial x} \right\rangle = -\left\langle \frac{\partial V}{\partial x} \right\rangle = \langle F(x) \rangle
\label{eq:ehrenfest_p}
\end{equation}

เมื่อรวมสมการ~\eqref{eq:ehrenfest_x} และ~\eqref{eq:ehrenfest_p}
\begin{equation}
m \frac{d^2}{dt^2}\langle x \rangle = \langle F(x) \rangle
\label{eq:quantum_newton}
\end{equation}

หากกระจายแรง $F(x) = -\frac{\partial V}{\partial x}$ ด้วยอนุกรมเทย์เลอร์รอบค่าคาดหวังตำแหน่ง $\langle x \rangle$
\begin{equation}
F(x) = F(\langle x \rangle) + F'(\langle x \rangle)(x - \langle x \rangle) + \frac{1}{2}F''(\langle x \rangle)(x - \langle x \rangle)^2 + \dots
\end{equation}
เมื่อหาค่าคาดหวัง พจน์ที่สองจะกลายเป็นศูนย์ และพจน์ที่สามขึ้นกับความแปรปรวน $\sigma_x^2$
\begin{equation}
\langle F(x) \rangle = F(\langle x \rangle) + \frac{1}{2} F''(\langle x \rangle) \sigma_x^2 + \dots
\end{equation}
ในกรณีที่กลุ่มคลื่นมีความกว้างแคบมาก ($\sigma_x \to 0$) หรือสนามศักย์เปลี่ยนแปลงช้ามากในระดับขนาดของกลุ่มคลื่น ($F'' \approx 0$) จะได้ $\langle F(x) \rangle \approx F(\langle x \rangle)$ ส่งผลให้สมการ~\eqref{eq:quantum_newton} ลดรูปกลับไปเป็น \textbf{กฎการเคลื่อนที่ข้อที่สองของนิวตัน} $m \frac{d^2 \langle x \rangle}{dt^2} = F(\langle x \rangle)$ อย่างแม่นยำ นี่คือการแสดงออกอย่างเป็นรูปธรรมของ \textbf{หลักความสอดคล้อง (Correspondence Principle)}

% ==========================================================
% 3.7 พรมแดนร่วมสมัย โครงข่ายประสาทเทียมฟิสิกส์ (PINNs)
% ==========================================================
\section{พรมแดนร่วมสมัย ---- การแก้สมการชเรอดิงเงอร์ด้วยโครงข่ายประสาทเทียมฟิสิกส์ (PINNs)}
\index{โครงข่ายประสาทเทียมฟิสิกส์ (Physics-Informed Neural Networks)}
\index{การหาอนุพันธ์อัตโนมัติ (Automatic Differentiation)}
\index{การเรียนรู้เชิงลึกทางควอนตัม (Quantum Deep Learning)}

\subsection{แนวคิดของ Physics-Informed Neural Networks (PINNs)}
ในการคำนวณทางฟิสิกส์ควอนตัมสมัยใหม่ ระเบียบวิธีดั้งเดิม เช่น วิธีผลต่างสืบเนื่อง (Finite Difference) หรือวิธีไฟไนต์เอลิเมนต์ (Finite Element) ประสบปัญหาการระเบิดเชิงมิติ (Curse of Dimensionality) เมื่อขยายไปสู่ระบบหลายอนุภาค 

ในคริสต์ศักราช 2019 มาซิอาร์ ไรส์ซี (Maziar Raissi) ปาริส เพอร์ดิคาริส (Paris Perdikaris) และ จอร์จ คาร์เนียดาคิส (George Em Karniadakis) ได้เสนอสถาปัตยกรรม \textbf{โครงข่ายประสาทเทียมฟิสิกส์ (Physics-Informed Neural Networks --- PINNs)} ซึ่งฝังกฎฟิสิกส์เชิงอนุพันธ์ย่อย (PDE) เข้าไปในฟังก์ชันการสูญเสียของโครงข่ายประสาทเทียมโดยตรงผ่านการหาอนุพันธ์อัตโนมัติ (Automatic Differentiation / Autograd)

\subsection{สถาปัตยกรรมและฟังก์ชันการสูญเสียสำหรับสมการชเรอดิงเงอร์}
ให้โครงข่ายประสาทเทียม $NN(x, t; \boldsymbol{\theta})$ ที่มีพารามิเตอร์เวตและไบอัส $\boldsymbol{\theta}$ ทำนายค่าฟังก์ชันคลื่นส่วนจริง $u(x, t)$ และส่วนจินตภาพ $v(x, t)$ โดยที่ $\Psi(x, t) = u(x, t) + i v(x, t)$

แยกสมการชเรอดิงเงอร์ขึ้นกับเวลา $i\hbar \frac{\partial \Psi}{\partial t} = -\frac{\hbar^2}{2m}\frac{\partial^2 \Psi}{\partial x^2} + V(x)\Psi$ ออกเป็นส่วนจริงและส่วนจินตภาพ
\begin{align}
f_u &= -\hbar \frac{\partial v}{\partial t} + \frac{\hbar^2}{2m}\frac{\partial^2 u}{\partial x^2} - V(x)u = 0 \\
f_v &= \hbar \frac{\partial u}{\partial t} + \frac{\hbar^2}{2m}\frac{\partial^2 v}{\partial x^2} - V(x)v = 0
\end{align}

ฟังก์ชันการสูญเสียรวม (Total Loss Function, $\mathcal{L}_{\text{total}}$) ประกอบด้วย 3 ส่วน
\begin{equation}
\mathcal{L}_{\text{total}}(\boldsymbol{\theta}) = w_{\text{pde}}\mathcal{L}_{\text{PDE}} + w_{\text{ic}}\mathcal{L}_{\text{IC}} + w_{\text{bc}}\mathcal{L}_{\text{BC}}
\label{eq:pinn_loss}
\end{equation}

โดยที่
\begin{enumerate}
    \item \textbf{การสูญเสียตามสมการฟิสิกส์ (PDE Residual Loss)} ---- บังคับให้โครงข่ายสอดคล้องกับสมการชเรอดิงเงอร์ ณ จุดจัดวางสุ่ม (Collocation Points, $N_f$)
    \begin{equation}
    \mathcal{L}_{\text{PDE}} = \frac{1}{N_f} \sum_{i=1}^{N_f} \left( |f_u(x_i, t_i)|^2 + |f_v(x_i, t_i)|^2 \right)
    \end{equation}
    \item \textbf{การสูญเสียเงื่อนไขเริ่มต้น (Initial Condition Loss)} ---- สอดคล้องกับสถานะเริ่มต้น $\Psi(x, 0) = \Psi_0(x)$ ณ จุดข้อมูล $N_0$ จุด
    \begin{equation}
    \mathcal{L}_{\text{IC}} = \frac{1}{N_0} \sum_{j=1}^{N_0} \left( |u(x_j, 0) - u_0(x_j)|^2 + |v(x_j, 0) - v_0(x_j)|^2 \right)
    \end{equation}
    \item \textbf{การสูญเสียเงื่อนไขขอบเขต (Boundary Condition Loss)} ---- บังคับเงื่อนไขขอบเขต เช่น $\Psi(\pm L, t) = 0$
    \begin{equation}
    \mathcal{L}_{\text{BC}} = \frac{1}{N_b} \sum_{k=1}^{N_b} \left( |\Psi(x_k^b, t_k)|^2 \right)
    \end{equation}
\end{enumerate}

การอัปเดตพารามิเตอร์ $\boldsymbol{\theta}$ ด้วยขั้นตอนวิธี Adam Optimizer ร่วมกับ L-BFGS ช่วยให้โครงข่ายประสาทเทียมสามารถเรียนรู้และค้นพบฟังก์ชันคลื่นควอนตัมได้โดยไม่ต้องสร้างกริดแบบดั้งเดิม (Mesh-free Solver) ซึ่งกำลังเป็นเครื่องมือแนวหน้าในการจำลองสสารควอนตัมหลายตัว (Many-Body Quantum Systems) และควอนตัมเคมีในปัจจุบัน

% ==========================================================
% กล่องปฏิบัติการ WebAR
% ==========================================================
\begin{qrcodelabbox}{การทดลองเสมือนจริง วิวัฒนาการเวลาของฟังก์ชันคลื่นควอนตัม WebAR}{https://tsanaphysics.github.io/ar-Quantum/lab03_schrodinger_wavepacket.html}
\begin{minipage}[c]{0.70\textwidth}
\textbf{คำแนะนำการทดลองเสมือนจริงสามมิติ (Interactive WebAR Simulation)}
\begin{enumerate}
    \item สแกน QR Code ด้วยอุปกรณ์สมาร์ทโฟนหรือแท็บเล็ตเพื่อเปิดแบบจำลอง WebAR
    \item สังเกตการแกว่งกวัดของฟังก์ชันคลื่นส่วนจริง ($u$), ส่วนจินตภาพ ($v$) และความหนาแน่นความน่าจะเป็น ($|\Psi|^2$) ใน 3 มิติ
    \item ทดลองปรับเปลี่ยนศักย์ขวางกั้น $V(x)$ เพื่อสังเกตการสะท้อนและการทะลุผ่านของกลุ่มคลื่น
    \item ตรวจสอบเวกเตอร์กระแสความน่าจะเป็น $\mathbf{j}(x, t)$ ที่เคลื่อนที่สอดคล้องกับสมการความต่อเนื่อง
\end{enumerate}
\end{minipage}
\hfill
\begin{minipage}[c]{0.25\textwidth}
\centering
\qrcode[height=2.8cm]{https://tsanaphysics.github.io/ar-Quantum/lab03_schrodinger_wavepacket.html}
\end{minipage}
\end{qrcodelabbox}

% ==========================================================
% ตัวอย่างโจทย์ P-S-C-T 7 ข้อเข้มข้น
% ==========================================================
\section{ตัวอย่างการแก้โจทย์เชิงวิเคราะห์ตามระเบียบวิธี P-S-C-T}

\begin{psctbox}{การวิเคราะห์สถานะซ้อนทับและคาบการแกว่งกวัดของความน่าจะเป็น}
\textbf{Problem (โจทย์ปัญหา)}

อนุภาคในบ่อศักย์ลึกอนันต์ 1 มิติกว้าง $L$ ($0 \le x \le L$) มีฟังก์ชันคลื่นเริ่มต้น ณ เวลา $t = 0$ อยู่ในสถานะซ้อนทับเชิงเส้นระหว่างสถานะพื้น ($n = 1$) และสถานะกระตุ้นแรก ($n = 2$) ดังนี้
\begin{equation*}
\Psi(x, 0) = \frac{1}{\sqrt{2}}\psi_1(x) + \frac{1}{\sqrt{2}}\psi_2(x)
\end{equation*}
โดยที่ $\psi_n(x) = \sqrt{\frac{2}{L}}\sin\left( \frac{n\pi x}{L} \right)$ และ $E_n = \frac{n^2 \pi^2 \hbar^2}{2mL^2}$
\begin{enumerate}
    \item จงหาฟังก์ชันคลื่น $\Psi(x, t)$ ณ เวลา $t > 0$
    \item จงคำนวณความหนาแน่นความน่าจะเป็น $|\Psi(x, t)|^2$ และแสดงว่ามีการแกว่งกวัดตามเวลา
    \item จงหาคาบเวลาการแกว่งกวัด ($T_{\text{osc}}$) และความถี่เชิงมุม ($\omega_{\text{osc}}$)
\end{enumerate}

\textbf{Solution (วิธีทำ)}

\textbf{ขั้นตอนที่ 1} ใส่แฟกเตอร์วิวัฒนาการเวลา $e^{-iE_n t/\hbar}$ ให้แก่แต่ละสถานะ
\begin{equation*}
\Psi(x, t) = \frac{1}{\sqrt{2}}\psi_1(x) e^{-iE_1 t/\hbar} + \frac{1}{\sqrt{2}}\psi_2(x) e^{-iE_2 t/\hbar}
\end{equation*}

\textbf{ขั้นตอนที่ 2} คำนวณความหนาแน่นความน่าจะเป็น $|\Psi(x, t)|^2 = \Psi^* \Psi$
\begin{align*}
|\Psi(x, t)|^2 &= \left( \frac{1}{\sqrt{2}}\psi_1 e^{iE_1 t/\hbar} + \frac{1}{\sqrt{2}}\psi_2 e^{iE_2 t/\hbar} \right) \left( \frac{1}{\sqrt{2}}\psi_1 e^{-iE_1 t/\hbar} + \frac{1}{\sqrt{2}}\psi_2 e^{-iE_2 t/\hbar} \right) \\
&= \frac{1}{2}\psi_1^2(x) + \frac{1}{2}\psi_2^2(x) + \frac{1}{2}\psi_1(x)\psi_2(x)\left( e^{i(E_2 - E_1)t/\hbar} + e^{-i(E_2 - E_1)t/\hbar} \right) \\
&= \frac{1}{2}\psi_1^2(x) + \frac{1}{2}\psi_2^2(x) + \psi_1(x)\psi_2(x)\cos\left( \frac{E_2 - E_1}{\hbar} t \right)
\end{align*}

\textbf{ขั้นตอนที่ 3} คำนวณความถี่และคาบของการแกว่งกวัด
ความถี่เชิงมุมของการแกว่งกวัดคือ
\begin{equation*}
\omega_{\text{osc}} = \frac{E_2 - E_1}{\hbar} = \frac{4E_1 - E_1}{\hbar} = \frac{3E_1}{\hbar} = \frac{3\pi^2 \hbar}{2mL^2}
\end{equation*}
คาบของการแกว่งกวัดคือ
\begin{equation*}
T_{\text{osc}} = \frac{2\pi}{\omega_{\text{osc}}} = \frac{2\pi \hbar}{3E_1} = \frac{4mL^2}{3\pi \hbar}
\end{equation*}

\textbf{Check (ตรวจสอบ)} ---- เมื่อ $t = 0$ พจน์ $\cos(0) = 1$ ทำให้ $|\Psi|^2 = \frac{1}{2}(\psi_1 + \psi_2)^2$ สอดคล้องกับสภาวะเริ่มต้น และความน่าจะเป็นรวม $\int_0^L |\Psi(x, t)|^2 dx = \frac{1}{2}(1) + \frac{1}{2}(1) + 0 = 1$ คงที่ตลอดเวลาเนื่องจาก $\psi_1$ และ $\psi_2$ ตั้งฉากกัน

\textbf{Think (คิดต่อ)} ---- ในสถานะคงตัวเดี่ยว ความหนาแน่นความน่าจะเป็นจะอยู่นิ่งไม่เปลี่ยนตามเวลา แต่เมื่อใดก็ตามที่มีการซ้อนทับของสถานะที่มีพลังงานต่างกัน ความน่าจะเป็นจะแกว่งกวัดไปมาระหว่างสองฟากของบ่อศักย์ด้วยความถี่บีต $\Delta E/\hbar$ ซึ่งเป็นกลไกพื้นฐานของการเกิดการแผ่รังสีแม่เหล็กไฟฟ้าจากการเปลี่ยนระดับพลังงานในอะตอม
\end{psctbox}

\begin{psctbox}{การคำนวณกระแสความน่าจะเป็นของคลื่นทรงกลมแผ่ออก}
\textbf{Problem (โจทย์ปัญหา)}

ในปรากฏการณ์การกระเจิงควอนตัม 3 มิติ ลำอนุภาคที่กระเจิงออกจากจุดศูนย์กลางสามารถแทนด้วยฟังก์ชันคลื่นทรงกลมแผ่ออก (Outgoing Spherical Wave)
\begin{equation*}
\Psi(r, t) = \frac{A}{r} e^{i(kr - \omega t)}
\end{equation*}
โดยที่ $A$ คือค่าคงตัวแอมพลิจูด และ $k = p/\hbar$ จงคำนวณ
\begin{enumerate}
    \item ความหนาแน่นกระแสความน่าจะเป็นในแนวรัศมี ($j_r$)
    \item ฟลักซ์ความน่าจะเป็นรวมทั้งหมดที่ไหลผ่านผิวทรงกลมรัศมี $R$ ปิดล้อมจุดกำเนิด
\end{enumerate}

\textbf{Solution (วิธีทำ)}

\textbf{ขั้นตอนที่ 1} ตัวดำเนินการเกรเดียนต์ในพิกัดทรงกลมสำหรับฟังก์ชันที่ขึ้นกับรัศมี $r$ เท่านั้นคือ $\nabla = \hat{r}\frac{\partial}{\partial r}$
หาอนุพันธ์ของฟังก์ชันคลื่น
\begin{equation*}
\frac{\partial \Psi}{\partial r} = \frac{\partial}{\partial r}\left( \frac{A}{r} e^{i(kr - \omega t)} \right) = \left( -\frac{A}{r^2} + \frac{ikA}{r} \right) e^{i(kr - \omega t)} = \left( ik - \frac{1}{r} \right)\Psi
\end{equation*}

\textbf{ขั้นตอนที่ 2} คำนวณความหนาแน่นกระแสความน่าจะเป็นจากสมการ~\eqref{eq:prob_current_3d}
\begin{align*}
\mathbf{j}(r) &= \frac{\hbar}{m}\text{Im}\left( \Psi^* \frac{\partial \Psi}{\partial r}\hat{r} \right) = \frac{\hbar}{m}\text{Im}\left[ \Psi^* \left( ik - \frac{1}{r} \right)\Psi \hat{r} \right] \\
&= \frac{\hbar}{m}\text{Im}\left[ |\Psi|^2 \left( ik - \frac{1}{r} \right) \right] \hat{r} = \frac{\hbar}{m} |\Psi|^2 k \, \hat{r}
\end{align*}
เนื่องจาก $|\Psi|^2 = \frac{|A|^2}{r^2}$ จะได้
\begin{equation*}
\mathbf{j}(r) = \frac{\hbar k}{m} \frac{|A|^2}{r^2} \hat{r} = v \frac{|A|^2}{r^2} \hat{r}
\end{equation*}
โดยที่ $v = \hbar k/m$ คือความเร็วของอนุภาค

\textbf{ขั้นตอนที่ 3} คำนวณฟลักซ์รวมผ่านผิวทรงกลมรัศมี $R$
\begin{equation*}
\Phi_{\text{prob}} = \oint_{S} \mathbf{j}\cdot d\mathbf{a} = \int_{0}^{2\pi}\int_{0}^{\pi} \left( v \frac{|A|^2}{R^2} \right) R^2 \sin\theta \, d\theta \, d\phi = 4\pi v |A|^2
\end{equation*}

\textbf{Check (ตรวจสอบ)} ---- ฟลักซ์รวม $\Phi_{\text{prob}} = 4\pi v |A|^2$ เป็นค่าคงที่ที่ไม่ขึ้นกับรัศมี $R$ เลย สอดคล้องกับกฎอนุรักษ์ความน่าจะเป็นอย่างสมบูรณ์ เพราะไม่มีความน่าจะเป็นสะสมหรือสูญหายในอวกาศระหว่างทาง

\textbf{Think (คิดต่อ)} ---- พจน์ $1/r^2$ ในกระแสความน่าจะเป็นสะท้อนกฎกำลังสองผกผันเชิงเรขาคณิต (Inverse-Square Law) ของความเข้มรังสี ซึ่งเป็นผลลัพธ์โดยตรงจากการอนุรักษ์ความน่าจะเป็นใน 3 มิติ
\end{psctbox}

\begin{psctbox}{พีชคณิตเมทริกซ์ในปริภูมิฮิลเบิร์ต 2 มิติ}
\textbf{Problem (โจทย์ปัญหา)}

ในปริภูมิฮิลเบิร์ต 2 มิติ กำหนดฐานตั้งฉากปรกติ $\{\ket{1}, \ket{2}\}$ ตัวดำเนินการ $\hat{\Omega}$ มีการแทนด้วยเมทริกซ์ดังนี้
\begin{equation*}
\boldsymbol{\Omega} = \begin{pmatrix} 2 & i \\ -i & 2 \end{pmatrix}
\end{equation*}
\begin{enumerate}
    \item จงแสดงว่า $\hat{\Omega}$ เป็นตัวดำเนินการแฮร์มิเชียน
    \item จงหาค่าเจาะจงทั้งหมด และเวกเตอร์สถานะเจาะจงที่ได้รับการนอร์มอลไลซ์
    \item จงตรวจสอบว่าสถานะเจาะจงทั้งสองตั้งฉากกันและสอดคล้องกับความสมบูรณ์ของฐาน
\end{enumerate}

\textbf{Solution (วิธีทำ)}

\textbf{ขั้นตอนที่ 1} ตรวจสอบความเป็นแฮร์มิเชียน $\boldsymbol{\Omega}^\dagger = (\boldsymbol{\Omega}^*)^T$
\begin{equation*}
\boldsymbol{\Omega}^* = \begin{pmatrix} 2 & -i \\ i & 2 \end{pmatrix} \implies (\boldsymbol{\Omega}^*)^T = \begin{pmatrix} 2 & i \\ -i & 2 \end{pmatrix} = \boldsymbol{\Omega}
\end{equation*}
ดังนั้น $\hat{\Omega}$ เป็นแฮร์มิเชียน

\textbf{ขั้นตอนที่ 2} หาค่าเจาะจงจากสมการลักษณะเฉพาะ $\det(\boldsymbol{\Omega} - \lambda \mathbf{I}) = 0$
\begin{equation*}
\begin{vmatrix} 2-\lambda & i \\ -i & 2-\lambda \end{vmatrix} = (2-\lambda)^2 - (i)(-i) = (2-\lambda)^2 - 1 = 0
\end{equation*}
\begin{equation*}
(2-\lambda)^2 = 1 \implies 2-\lambda = \pm 1 \implies \lambda_1 = 3, \quad \lambda_2 = 1
\end{equation*}
ค่าเจาะจงทั้งสองเป็นจำนวนจริง สอดคล้องกับทฤษฎีบทของแฮร์มิเชียน

\textbf{ขั้นตอนที่ 3} หาเวกเตอร์เจาะจงสำหรับ $\lambda_1 = 3$
\begin{equation*}
\begin{pmatrix} 2-3 & i \\ -i & 2-3 \end{pmatrix} \begin{pmatrix} u_1 \\ u_2 \end{pmatrix} = \begin{pmatrix} -1 & i \\ -i & -1 \end{pmatrix} \begin{pmatrix} u_1 \\ u_2 \end{pmatrix} = \begin{pmatrix} 0 \\ 0 \end{pmatrix} \implies -u_1 + i u_2 = 0 \implies u_1 = i u_2
\end{equation*}
นอร์มอลไลซ์ $|u_1|^2 + |u_2|^2 = |i u_2|^2 + |u_2|^2 = 2|u_2|^2 = 1 \implies u_2 = \frac{1}{\sqrt{2}}$
\begin{equation*}
\ket{\lambda_1} = \frac{1}{\sqrt{2}}\begin{pmatrix} i \\ 1 \end{pmatrix}
\end{equation*}

หาเวกเตอร์เจาะจงสำหรับ $\lambda_2 = 1$
\begin{equation*}
\begin{pmatrix} 2-1 & i \\ -i & 2-1 \end{pmatrix} \begin{pmatrix} v_1 \\ v_2 \end{pmatrix} = \begin{pmatrix} 1 & i \\ -i & 1 \end{pmatrix} \begin{pmatrix} v_1 \\ v_2 \end{pmatrix} = \begin{pmatrix} 0 \\ 0 \end{pmatrix} \implies v_1 + i v_2 = 0 \implies v_1 = -i v_2
\end{equation*}
\begin{equation*}
\ket{\lambda_2} = \frac{1}{\sqrt{2}}\begin{pmatrix} -i \\ 1 \end{pmatrix}
\end{equation*}

\textbf{ขั้นตอนที่ 4} ตรวจสอบความตั้งฉากและความสมบูรณ์
\begin{equation*}
\braket{\lambda_1}{\lambda_2} = \frac{1}{\sqrt{2}}\begin{pmatrix} -i & 1 \end{pmatrix} \frac{1}{\sqrt{2}}\begin{pmatrix} -i \\ 1 \end{pmatrix} = \frac{1}{2}\big( (-i)(-i) + (1)(1) \big) = \frac{1}{2}(-1 + 1) = 0
\end{equation*}
และตัวดำเนินการฉายภาพรวม
\begin{align*}
\ket{\lambda_1}\bra{\lambda_1} + \ket{\lambda_2}\bra{\lambda_2} &= \frac{1}{2}\begin{pmatrix} i \\ 1 \end{pmatrix}\begin{pmatrix} -i & 1 \end{pmatrix} + \frac{1}{2}\begin{pmatrix} -i \\ 1 \end{pmatrix}\begin{pmatrix} i & 1 \end{pmatrix} \\
&= \frac{1}{2}\begin{pmatrix} 1 & i \\ -i & 1 \end{pmatrix} + \frac{1}{2}\begin{pmatrix} 1 & -i \\ i & 1 \end{pmatrix} = \begin{pmatrix} 1 & 0 \\ 0 & 1 \end{pmatrix} = \mathbf{I}
\end{align*}

\textbf{Check (ตรวจสอบ)} ---- ค่าเจาะจงเป็นจริง เวกเตอร์เจาะจงตั้งฉากกัน และผลบวกของโปรเจกเตอร์เท่ากับเมทริกซ์เอกลักษณ์อย่างแม่นยำ

\textbf{Think (คิดต่อ)} ---- เมทริกซ์ลักษณะนี้คือรูปแบบทั่วไปของการแทนคิวบิตในคอมพิวเตอร์ควอนตัม ซึ่งสามารถแยกออกเป็นส่วนผสมของเมทริกซ์พอลิ (Pauli Matrices) $\boldsymbol{\Omega} = 2\mathbf{I} - \sigma_y$
\end{psctbox}

\begin{psctbox}{การประยุกต์ทฤษฎีบทเอห์เรนเฟสต์กับฮาร์มอนิกออสซิลเลเตอร์}
\textbf{Problem (โจทย์ปัญหา)}

สำหรับตัวแกว่งกวัดฮาร์มอนิก 1 มิติ ฮามิลโทเนียนกำหนดโดย $\hat{H} = \frac{\hat{p}^2}{2m} + \frac{1}{2}m\omega^2 \hat{x}^2$
\begin{enumerate}
    \item จงใช้ทฤษฎีบทเอห์เรนเฟสต์เพื่อสร้างสมการเชิงอนุพันธ์ของ $\frac{d}{dt}\langle x \rangle$ และ $\frac{d}{dt}\langle p \rangle$
    \item จงหาผลเฉลยของ $\langle x(t) \rangle$ สำหรับสภาวะเริ่มต้นใดๆ ที่มี $\langle x(0) \rangle = x_0$ และ $\langle p(0) \rangle = p_0$
\end{enumerate}

\textbf{Solution (วิธีทำ)}

\textbf{ขั้นตอนที่ 1} คำนวณคอมมิวเตเตอร์
จากสมการ~\eqref{eq:ehrenfest_x} ทราบว่า $\frac{d}{dt}\langle x \rangle = \frac{\langle p \rangle}{m}$
สำหรับโมเมนตัม คำนวณคอมมิวเตเตอร์กับพลังงานศักย์ $V(x) = \frac{1}{2}m\omega^2 x^2$
\begin{equation*}
[\hat{p}, \hat{H}] = \left[ \hat{p}, \frac{1}{2}m\omega^2 \hat{x}^2 \right] = \frac{1}{2}m\omega^2 [\hat{p}, \hat{x}^2] = \frac{1}{2}m\omega^2 (-2i\hbar \hat{x}) = -i\hbar m\omega^2 \hat{x}
\end{equation*}
ใช้สมการเอห์เรนเฟสต์
\begin{equation*}
\frac{d}{dt}\langle p \rangle = \frac{1}{i\hbar}\langle [\hat{p}, \hat{H}] \rangle = -m\omega^2 \langle x \rangle
\end{equation*}

\textbf{ขั้นตอนที่ 2} แก้ระบบสมการอนุพันธ์ควบคู่
หาอนุพันธ์เทียบกับเวลาอีกหนึ่งครั้งของสมการแรก
\begin{equation*}
\frac{d^2}{dt^2}\langle x \rangle = \frac{1}{m}\frac{d}{dt}\langle p \rangle = \frac{1}{m}(-m\omega^2 \langle x \rangle) = -\omega^2 \langle x \rangle
\end{equation*}
นี่คือสมการการแกว่งกวัดฮาร์มอนิกอย่างง่าย ผลเฉลยทั่วไปคือ
\begin{equation*}
\langle x(t) \rangle = A \cos(\omega t) + B \sin(\omega t)
\end{equation*}

\textbf{ขั้นตอนที่ 3} กำหนดค่าคงตัวจากเงื่อนไขเริ่มต้น
ที่ $t = 0$ จะได้ $\langle x(0) \rangle = A = x_0$
หาอนุพันธ์เพื่อเชื่อมโยงกับโมเมนตัม
\begin{equation*}
\frac{d}{dt}\langle x(0) \rangle = \omega B = \frac{\langle p(0) \rangle}{m} = \frac{p_0}{m} \implies B = \frac{p_0}{m\omega}
\end{equation*}
ดังนั้น ผลเฉลยสมบูรณ์คือ
\begin{equation*}
\langle x(t) \rangle = x_0 \cos(\omega t) + \frac{p_0}{m\omega}\sin(\omega t)
\end{equation*}

\textbf{Check (ตรวจสอบ)} ---- ผลเฉลยของค่าคาดหวังตำแหน่ง $\langle x(t) \rangle$ มีรูปแบบเหมือนกับสมการการเคลื่อนที่ของมวลติดสปริงในกลศาสตร์คลาสสิกทุกประการ

\textbf{Think (คิดต่อ)} ---- ในกรณีของศักย์ฮาร์มอนิก อนุพันธ์อันดับสองของแรง $F''(x) = -\frac{d^3 V}{dx^3} = 0$ ทำให้พจน์ความคลาดเคลื่อนทางควอนตัมเป็นศูนย์อย่างแท้จริง ค่าคาดหวังของฮาร์มอนิกออสซิลเลเตอร์จึงเคลื่อนที่ตามวิถีคลาสสิกได้อย่างสมบูรณ์แบบไม่ว่ากลุ่มคลื่นจะกว้างเพียงใด
\end{psctbox}

\begin{psctbox}{การแปลงฟูเรียร์สู่ปริภูมิโมเมนตัมของอนุภาคในกล่อง}
\textbf{Problem (โจทย์ปัญหา)}

อนุภาคในกล่อง 1 มิติที่สถานะพื้น ($n = 1$) มีฟังก์ชันคลื่นในปริภูมิตำแหน่งคือ
\begin{equation*}
\psi(x) = \begin{cases} \sqrt{\frac{2}{L}}\sin\left( \frac{\pi x}{L} \right) & \text{เมื่อ } 0 \le x \le L \\ 0 & \text{กรณีอื่น} \end{cases}
\end{equation*}
จงหาฟังก์ชันคลื่นในปริภูมิโมเมนตัม $\Phi(p)$ และคำนวณการกระจายความน่าจะเป็นที่จะวัดได้โมเมนตัม $p$

\textbf{Solution (วิธีทำ)}

\textbf{ขั้นตอนที่ 1} ฟังก์ชันคลื่นในปริภูมิโมเมนตัมหาได้จากการฉายภาพ $\Phi(p) = \braket{p}{\psi}$
\begin{equation*}
\Phi(p) = \frac{1}{\sqrt{2\pi\hbar}} \int_{-\infty}^{\infty} \psi(x) e^{-ipx/\hbar} \, dx = \frac{1}{\sqrt{2\pi\hbar}} \sqrt{\frac{2}{L}} \int_{0}^{L} \sin\left( \frac{\pi x}{L} \right) e^{-ipx/\hbar} \, dx
\end{equation*}

\textbf{ขั้นตอนที่ 2} ใช้เอกลักษณ์ออยเลอร์ $\sin(\theta) = \frac{e^{i\theta} - e^{-i\theta}}{2i}$
\begin{equation*}
\Phi(p) = \frac{1}{\sqrt{\pi\hbar L}}\frac{1}{2i} \int_{0}^{L} \left( e^{i(\pi/L - p/\hbar)x} - e^{-i(\pi/L + p/\hbar)x} \right) dx
\end{equation*}
อินทิเกรตแต่ละพจน์และแทนขอบเขตจาก $0$ ถึง $L$ หลังจัดรูปพีชคณิต
\begin{equation*}
\Phi(p) = \sqrt{\frac{\pi L}{\hbar}} \left( \frac{1 + e^{-ipL/\hbar}}{\pi^2 - (pL/\hbar)^2} \right)
\end{equation*}

\textbf{ขั้นตอนที่ 3} หาความหนาแน่นความน่าจะเป็นในปริภูมิโมเมนตัม $P(p) = |\Phi(p)|^2$
\begin{equation*}
|1 + e^{-ipL/\hbar}|^2 = (1 + \cos(pL/\hbar))^2 + \sin^2(pL/\hbar) = 2 + 2\cos(pL/\hbar) = 4\cos^2\left( \frac{pL}{2\hbar} \right)
\end{equation*}
ดังนั้น
\begin{equation*}
P(p) = \frac{4\pi L}{\hbar} \frac{\cos^2(pL / 2\hbar)}{\left[ \pi^2 - (pL/\hbar)^2 \right]^2}
\end{equation*}

\textbf{Check (ตรวจสอบ)} ---- เมื่อ $p \to \pm \frac{\pi\hbar}{L}$ (โมเมนตัมตามแบบจำลองโบร์) ใช้กฎโลปิตาลจะได้ค่าความน่าจะเป็นมีจุดยอดสูงสุด และเมื่ออินทิเกรตตลอดทุกค่าของโมเมนตัม $\int_{-\infty}^\infty P(p) dp = 1$ สอดคล้องกับทฤษฎีบทพาร์เซวาล (Parseval's Theorem)

\textbf{Think (คิดต่อ)} ---- แม้ในกลศาสตร์คลาสสิกอนุภาคที่สะท้อนไปมาจะมีโมเมนตัมเพียงสองค่าคือ $+p$ และ $-p$ แต่ในกลศาสตร์ควอนตัม เนื่องจากการกักขังอนุภาคในบริเวณจำกัด $L$ ทำให้เกิดการเลี้ยวเบนของคลื่น ส่งผลให้การกระจายโมเมนตัมทอดยาวออกไปทุกช่วงค่าตามหลักความไม่แน่นอนของไฮเซนเบิร์ก
\end{psctbox}

\begin{psctbox}{กระบวนการสร้างฐานตั้งฉากปรกติแบบแกรม-ชมิดต์สำหรับสถานะลดหลั่น}
\textbf{Problem (โจทย์ปัญหา)}

ในระบบควอนตัม 3 มิติ ฮามิลโทเนียนมีระดับพลังงานหนึ่งที่มีความลดหลั่น (Degeneracy) 2 เท่า โดยมีสถานะเจาะจงที่ไม่อิสระเชิงเส้นต่อกันสองสถานะ ได้แก่ $\ket{\phi_1}$ และ $\ket{\phi_2}$ ซึ่งได้รับการนอร์มอลไลซ์แล้ว ($\braket{\phi_1}{\phi_1} = \braket{\phi_2}{\phi_2} = 1$) แต่ไม่ได้ตั้งฉากกัน โดยมีผลคูณภายใน $\braket{\phi_1}{\phi_2} = \gamma = \frac{1}{2}$
จงใช้กระบวนการแกรม-ชมิดต์ (Gram-Schmidt Orthogonalization) เพื่อสร้างเซตฐานตั้งฉากปรกติ $\{\ket{\psi_1}, \ket{\psi_2}\}$

\textbf{Solution (วิธีทำ)}

\textbf{ขั้นตอนที่ 1} กำหนดให้สถานะแรกเป็นสถานะฐานตัวแรก
\begin{equation*}
\ket{\psi_1} = \ket{\phi_1}
\end{equation*}

\textbf{ขั้นตอนที่ 2} สร้างสถานะที่สองโดยการลบส่วนที่ฉายลงบน $\ket{\psi_1}$ ออกจาก $\ket{\phi_2}$
\begin{equation*}
\ket{\psi_2'} = \ket{\phi_2} - \braket{\psi_1}{\phi_2}\ket{\psi_1} = \ket{\phi_2} - \gamma \ket{\phi_1} = \ket{\phi_2} - \frac{1}{2}\ket{\phi_1}
\end{equation*}

\textbf{ขั้นตอนที่ 3} ตรวจสอบความตั้งฉากระหว่าง $\ket{\psi_1}$ และ $\ket{\psi_2'}$
\begin{equation*}
\braket{\psi_1}{\psi_2'} = \braket{\phi_1}{\phi_2 - \frac{1}{2}\phi_1} = \braket{\phi_1}{\phi_2} - \frac{1}{2}\braket{\phi_1}{\phi_1} = \frac{1}{2} - \frac{1}{2}(1) = 0
\end{equation*}
ตั้งฉากกันเรียบร้อยแล้ว

\textbf{ขั้นตอนที่ 4} นอร์มอลไลซ์ $\ket{\psi_2'}$ เพื่อหา $\ket{\psi_2}$
\begin{align*}
\braket{\psi_2'}{\psi_2'} &= \left( \bra{\phi_2} - \frac{1}{2}\bra{\phi_1} \right) \left( \ket{\phi_2} - \frac{1}{2}\ket{\phi_1} \right) \\
&= \braket{\phi_2}{\phi_2} - \frac{1}{2}\braket{\phi_2}{\phi_1} - \frac{1}{2}\braket{\phi_1}{\phi_2} + \frac{1}{4}\braket{\phi_1}{\phi_1} \\
&= 1 - \frac{1}{2}\left(\frac{1}{2}\right) - \frac{1}{2}\left(\frac{1}{2}\right) + \frac{1}{4}(1) = 1 - \frac{1}{4} - \frac{1}{4} + \frac{1}{4} = \frac{3}{4}
\end{align*}
ดังนั้น ตัวประกอบการนอร์มอลไลซ์คือ $\frac{1}{\sqrt{3/4}} = \frac{2}{\sqrt{3}}$
\begin{equation*}
\ket{\psi_2} = \frac{\ket{\psi_2'}}{\sqrt{\braket{\psi_2'}{\psi_2'}}} = \frac{2}{\sqrt{3}}\left( \ket{\phi_2} - \frac{1}{2}\ket{\phi_1} \right) = \frac{2}{\sqrt{3}}\ket{\phi_2} - \frac{1}{\sqrt{3}}\ket{\phi_1}
\end{equation*}

\textbf{Check (ตรวจสอบ)} ---- คำนวณนอร์ม $\braket{\psi_2}{\psi_2} = \frac{4}{3}(1) + \frac{1}{3}(1) - \frac{4}{3}(\frac{1}{2}) = \frac{5}{3} - \frac{2}{3} = 1$ และ $\braket{\psi_1}{\psi_2} = \frac{2}{\sqrt{3}}(\frac{1}{2}) - \frac{1}{\sqrt{3}}(1) = 0$ สมบูรณ์

\textbf{Think (คิดต่อ)} ---- ในทฤษฎีการรบกวนสำหรับสถานะลดหลั่น (Degenerate Perturbation Theory) การเลือกฐานตั้งฉากปรกติที่ถูกต้องที่ทำให้ตัวดำเนินการรบกวนอยู่ในรูปเมทริกซ์ทแยงมุมมีความสำคัญอย่างยิ่งยวดต่อการหาการแยกของระดับพลังงาน
\end{psctbox}

\begin{psctbox}{การคำนวณฟังก์ชันการสูญเสียของ PINNs สำหรับสมการชเรอดิงเงอร์}
\textbf{Problem (โจทย์ปัญหา)}

ในการฝึกโครงข่ายประสาทเทียมฟิสิกส์ (PINNs) เพื่อแก้สมการชเรอดิงเงอร์ของอนุภาคอิสระหนึ่งมิติ ($m = 1, \hbar = 1, V = 0$) สมมติว่าโครงข่ายประสาทเทียมทำนายค่าฟังก์ชันคลื่นทดสอบ ณ จุดจัดวางหนึ่ง $x = 1.0, t = 0.5$ ได้ค่าอนุพันธ์ย่อยจากการหาอนุพันธ์อัตโนมัติ (Autograd) ดังนี้
\begin{equation*}
\frac{\partial u}{\partial t} = 0.40, \quad \frac{\partial v}{\partial t} = -0.60, \quad \frac{\partial^2 u}{\partial x^2} = -1.20, \quad \frac{\partial^2 v}{\partial x^2} = 0.80
\end{equation*}
\begin{enumerate}
    \item จงคำนวณเศษตกค้างของสมการฟิสิกส์ $f_u$ และ $f_v$ ณ จุดนี้
    \item จงคำนวณค่าการสูญเสียฟิสิกส์เฉพาะจุด $\mathcal{L}_{\text{PDE}}(x, t) = |f_u|^2 + |f_v|^2$
\end{enumerate}

\textbf{Solution (วิธีทำ)}

\textbf{ขั้นตอนที่ 1} จากสมการเศษตกค้างของ TDSE สำหรับอนุภาคอิสระ ($m = 1, \hbar = 1, V = 0$)
\begin{align*}
f_u &= -\frac{\partial v}{\partial t} + \frac{1}{2}\frac{\partial^2 u}{\partial x^2} \\
f_v &= \frac{\partial u}{\partial t} + \frac{1}{2}\frac{\partial^2 v}{\partial x^2}
\end{align*}

\textbf{ขั้นตอนที่ 2} แทนค่าตัวเลขของอนุพันธ์ย่อย
\begin{equation*}
f_u = -(-0.60) + \frac{1}{2}(-1.20) = 0.60 - 0.60 = 0.00
\end{equation*}
\begin{equation*}
f_v = 0.40 + \frac{1}{2}(0.80) = 0.40 + 0.40 = 0.80
\end{equation*}

\textbf{ขั้นตอนที่ 3} คำนวณค่าความสูญเสียเฉพาะจุด
\begin{equation*}
\mathcal{L}_{\text{PDE}}(1.0, 0.5) = |f_u|^2 + |f_v|^2 = (0.00)^2 + (0.80)^2 = 0.64
\end{equation*}

\textbf{Check (ตรวจสอบ)} ---- ค่า $f_u = 0$ แสดงว่าส่วนจริงของสมการชเรอดิงเงอร์สอดคล้องกับฟิสิกส์พอดี แต่ $f_v = 0.80$ แสดงว่าส่วนจินตภาพยังมีค่าคลาดเคลื่อน จึงทำให้เกิดค่าความสูญเสีย $0.64$ ส่งสัญญาณให้เกิดการแพร่ย้อนกลับ (Backpropagation)

\textbf{Think (คิดต่อ)} ---- ในกระบวนการเรียนรู้จริง โครงข่ายจะสุ่มตัวอย่างจุด Collocation นับหมื่นจุด และใช้วิธีเกรเดียนต์เพื่อปรับค่าน้ำหนักจนกระทั่งค่าเฉลี่ยของ $\mathcal{L}_{\text{PDE}} < 10^{-5}$ ซึ่งจะทำให้ฟังก์ชันคลื่นที่โครงข่ายทำนายมีความถูกต้องแม่นยำเทียบเท่ากับผลเฉลยเชิงวิเคราะห์
\end{psctbox}

% ==========================================================
% 3.9 สรุปสาระสำคัญประจำบท
% ==========================================================
\section{สรุปสาระสำคัญประจำบทที่ 3}

\begin{enumerate}
    \item \textbf{สมการชเรอดิงเงอร์ขึ้นกับเวลา (TDSE)} ---- $i\hbar \frac{\partial \Psi}{\partial t} = \hat{H}\Psi$ เป็นสมการคลื่นเชิงเส้นที่ควบคุมพลวัตของสถานะควอนตัมทั้งหมด ความเป็นเชิงเส้นของสมการนำไปสู่หลักการซ้อนทับควอนตัม
    \item \textbf{สมการชเรอดิงเงอร์ไม่ขึ้นกับเวลา (TISE)} ---- การแยกตัวแปรสำหรับศักย์คงตัว $\Psi(x, t) = \psi(x)e^{-iEt/\hbar}$ นำไปสู่สมการค่าเจาะจง $\hat{H}\psi = E\psi$ ผลเฉลยคือสถานะคงตัวซึ่งมีความหนาแน่นความน่าจะเป็นคงที่ในเวลา
    \item \textbf{กฎอนุรักษ์ความน่าจะเป็น} ---- ความหนาแน่นความน่าจะเป็น $\rho = |\Psi|^2$ และกระแสความน่าจะเป็น $\mathbf{j} = \frac{\hbar}{m}\text{Im}(\Psi^*\nabla\Psi)$ สอดคล้องกับสมการความต่อเนื่อง $\frac{\partial\rho}{\partial t} + \nabla\cdot\mathbf{j} = 0$ ซึ่งรับประกันว่าผลรวมความน่าจะเป็นทั่วเอกภพเท่ากับ 1 เสมอตลอดกาล
    \item \textbf{สัญกรณ์ดิแรกและปริภูมิฮิลเบิร์ต} ---- เวกเตอร์สถานะแทนด้วยเคต $\ket{\psi}$ และดูอัลเวกเตอร์คือบรา $\bra{\psi}$ ผลคูณภายใน $\braket{\phi}{\psi}$ นิยามเรขาคณิตของปริภูมิ และฟังก์ชันคลื่นคือการฉายภาพ $\psi(x) = \braket{x}{\psi}$
    \item \textbf{ตัวดำเนินการแฮร์มิเชียน} ---- ปริมาณที่สังเกตวัดได้ทางกายภาพทั้งหมดแทนด้วยตัวดำเนินการแฮร์มิเชียน $\hat{A} = \hat{A}^\dagger$ ซึ่งให้ค่าเจาะจงเป็นจำนวนจริงเสมอ และสถานะเจาะจงที่ค่าเจาะจงต่างกันจะตั้งฉากกันอย่างสมบูรณ์
    \item \textbf{ทฤษฎีบทเอห์เรนเฟสต์} ---- อัตราการเปลี่ยนแปลงตามเวลาของค่าคาดหวังเป็นไปตาม $\frac{d}{dt}\langle A \rangle = \frac{1}{i\hbar}\langle [\hat{A}, \hat{H}] \rangle + \langle \frac{\partial \hat{A}}{\partial t} \rangle$ แสดงให้เห็นว่าค่าเฉลี่ยของตัวแปรควอนตัมเคลื่อนที่สอดคล้องกับกฎของนิวตัน
    \item \textbf{โครงข่ายประสาทเทียมฟิสิกส์ (PINNs)} ---- การปฏิวัติการแก้สมการชเรอดิงเงอร์เชิงตัวเลขด้วยการเรียนรู้เชิงลึก โดยการฝังสมการ PDE และเงื่อนไขขอบเขตลงในฟังก์ชันการสูญเสียของโครงข่ายประสาทเทียม
\end{enumerate}

% ==========================================================
% คำถามเชิงมโนทัศน์ 20 ข้อ
% ==========================================================
\section{คำถามเชิงมโนทัศน์ประจำบทที่ 3}

\begin{enumerate}
    \item เหตุใดสมการชเรอดิงเงอร์จึงต้องมีจำนวนจินตภาพ $i$ ปรากฏอยู่ในสมการ และจะเกิดผลอย่างไรต่อฟิสิกส์หากตัด $i$ ออกไป
    \item จงอธิบายความแตกต่างระหว่างสถานะคงตัว (Stationary State) และสถานะที่ไม่คงตัว ในแง่ของความหนาแน่นความน่าจะเป็นและค่าคาดหวัง
    \item เหตุใดฟังก์ชันคลื่น $\Psi(x, t)$ จึงไม่สามารถวัดได้โดยตรงในห้องปฏิบัติการ แต่ปริมาณ $|\Psi|^2$ จึงสามารถวัดได้
    \item ในสมการความต่อเนื่องของกระแสความน่าจะเป็น ความหมายทางกายภาพของพจน์ $\nabla\cdot\mathbf{j}$ คืออะไร
    \item หากเพิ่มค่าคงที่ $V_0$ เข้าไปในพลังงานศักย์ทั่วทั้งระบบ ($V'(x) = V(x) + V_0$) ฟังก์ชันคลื่นและค่าพลังงานจะเปลี่ยนแปลงไปอย่างไร
    \item จงอธิบายว่าเหตุใดตัวดำเนินการที่แทนปริมาณวัดได้ทางฟิสิกส์จึงต้องเป็นตัวดำเนินการแฮร์มิเชียนเท่านั้น
    \item ในสัญกรณ์ดิแรก ผลคูณภายใน $\braket{\phi}{\psi}$ และผลคูณภายนอก $\ket{\psi}\bra{\phi}$ มีความหมายทางเรขาคณิตและพีชคณิตแตกต่างกันอย่างไร
    \item ความสัมพันธ์ความสมบูรณ์ $\sum_n \ket{n}\bra{n} = \hat{I}$ มีประโยชน์อย่างไรในการเปลี่ยนฐานและการคำนวณในปริภูมิฮิลเบิร์ต
    \item จงอธิบายความหมายทางฟิสิกส์ของค่าคอมมิวเตเตอร์ $[\hat{A}, \hat{B}] = 0$ เทียบกับ $[\hat{A}, \hat{B}] \ne 0$ ในแง่ของการวัดพร้อมกัน
    \item ทฤษฎีบทเอห์เรนเฟสต์ยืนยันว่ากฎของนิวตันเป็นจริงเสมอในระดับควอนตัมหรือไม่ มีเงื่อนไขหรือข้อยกเว้นอย่างไร
    \item สำหรับอนุภาคอิสระ เหตุใดคลื่นระนาบ $e^{i(kx-\omega t)}$ จึงไม่สามารถนอร์มอลไลซ์ได้ในปริภูมิฮิลเบิร์ตแบบปกติ และนักฟิสิกส์แก้ปัญหานี้อย่างไร
    \item ในการแปลงเกจของสนามแม่เหล็กไฟฟ้า เหตุใดฟังก์ชันคลื่นจึงต้องเปลี่ยนเฟสเพื่อรักษากฎฟิสิกส์ให้คงเดิม
    \item หากสถานะควอนตัมอยู่ในสถานะเจาะจงของพลังงาน การวัดพลังงานซ้ำหลายๆ ครั้งจะได้ผลลัพธ์เป็นอย่างไร
    \item จงอธิบายบทบาทของตัวดำเนินการฉายภาพ $\hat{P} = \ket{\psi}\bra{\psi}$ ในกระบวนการวัดควอนตัมตามการแปลความหมายแบบโคเปนเฮเกน
    \item เหตุใดความหนาแน่นกระแสความน่าจะเป็นของคลื่นนิ่งจึงมีค่าเป็นศูนย์ตลอดเวลา แม้ว่าฟังก์ชันคลื่นจะมีการสั่นตามเวลาก็ตาม
    \item จงอธิบายความแตกต่างระหว่างสเปกตรัมแบบไม่ต่อเนื่อง (Discrete Spectrum) และสเปกตรัมแบบต่อเนื่อง (Continuous Spectrum) ของตัวดำเนินการ
    \item โครงข่ายประสาทเทียมฟิสิกส์ (PINNs) แตกต่างจากการเรียนรู้ของเครื่องแบบดั้งเดิมที่ใช้ข้อมูลฝึกสอน (Data-driven ML) อย่างไร
    \item การหาอนุพันธ์อัตโนมัติ (Autograd) ใน PINNs มีข้อได้เปรียบกว่าการหาอนุพันธ์แบบผลต่างสืบเนื่อง (Finite Difference) อย่างไร
    \item เหตุใดการฝึก PINNs สำหรับสมการชเรอดิงเงอร์จึงต้องแยกฟังก์ชันคลื่นออกเป็นส่วนจริงและส่วนจินตภาพ
    \item ในระบบหลายมิติ เหตุใด PINNs จึงสามารถก้าวข้ามข้อจำกัดการระเบิดเชิงมิติ (Curse of Dimensionality) ได้ดีกว่าวิธีเชิงตัวเลขแบบกริด
\end{enumerate}

% ==========================================================
% แบบฝึกหัดคำนวณและพิสูจน์เชิงวิเคราะห์ 20 ข้อ
% ==========================================================
\section{แบบฝึกหัดคำนวณและพิสูจน์เชิงวิเคราะห์ประจำบทที่ 3}

\begin{enumerate}
    \item อนุภาคตัวหนึ่งมีฟังก์ชันคลื่น $\psi(x) = A e^{-\lambda |x|}$ โดยที่ $A$ และ $\lambda$ เป็นค่าคงที่จริงบวก จงหาค่า $A$ ที่ทำให้นอร์มอลไลซ์ และคำนวณความน่าจะเป็นที่จะพบอนุภาคในช่วง $-\frac{1}{\lambda} \le x \le \frac{1}{\lambda}$
    \item จงคำนวณหาความหนาแน่นกระแสความน่าจะเป็น $j(x)$ ของฟังก์ชันคลื่น $\psi(x) = A e^{i k x} + B e^{-i k x}$ เมื่อ $A$ และ $B$ เป็นจำนวนเชิงซ้อน
    \item จงพิสูจน์ว่าตัวดำเนินการโมเมนตัม $\hat{p}_x = -i\hbar\frac{\partial}{\partial x}$ เป็นตัวดำเนินการแฮร์มิเชียนในปริภูมิของฟังก์ชันคลื่นที่ลู่เข้าสู่ศูนย์ที่ระยะอนันต์
    \item จงคำนวณคอมมิวเตเตอร์ $[\hat{x}^3, \hat{p}_x]$ และ $[\hat{x}, \hat{p}_x^2]$ โดยใช้เอกลักษณ์พีชคณิตคอมมิวเตเตอร์
    \item อนุภาคมีฟังก์ชันคลื่น $\Psi(x, t) = \sqrt{\frac{a}{\pi}} \frac{e^{-iEt/\hbar}}{x^2 + a^2}$ จงคำนวณหาค่าคาดหวัง $\langle x \rangle$ และ $\langle x^2 \rangle$
    \item จงแสดงว่า หาก $\hat{A}$ เป็นตัวดำเนินการแฮร์มิเชียน แล้ว ตัวดำเนินการ $\hat{U} = e^{i\hat{A}}$ จะเป็นตัวดำเนินการยูนิแทรี ($\hat{U}^\dagger \hat{U} = \hat{I}$)
    \item กำหนดเมทริกซ์ฮามิลโทเนียน $\mathbf{H} = \begin{pmatrix} E_0 & V \\ V & E_0 \end{pmatrix}$ จงหาค่าเจาะจงพลังงานและสถานะเจาะจงที่ได้รับการนอร์มอลไลซ์
    \item ในระบบ 2 มิติ สถานะเริ่มต้นคือ $\ket{\Psi(0)} = \begin{pmatrix} 1 \\ 0 \end{pmatrix}$ ภายใต้ฮามิลโทเนียนในข้อ 7 จงหาเคตสถานะ $\ket{\Psi(t)}$ ณ เวลา $t > 0$
    \item จงพิสูจน์เอกลักษณ์แฮดามาร์ด (Hadamard Lemma) $e^{\hat{A}}\hat{B}e^{-\hat{A}} = \hat{B} + [\hat{A}, \hat{B}] + \frac{1}{2!}[\hat{A}, [\hat{A}, \hat{B}]] + \dots$
    \item จงคำนวณกระแสความน่าจะเป็นสำหรับฟังก์ชันคลื่นเกาส์เซียน $\psi(x) = \left(\frac{1}{2\pi\sigma^2}\right)^{1/4} e^{-x^2/(4\sigma^2)} e^{ik_0 x}$
    \item อนุภาคอยู่ภายใต้ศักย์สนามโน้มถ่วงสม่ำเสมอ $V(x) = mgx$ จงใช้ทฤษฎีบทเอห์เรนเฟสต์พิสูจน์ว่า $\frac{d^2}{dt^2}\langle x \rangle = -g$
    \item จงพิสูจน์ว่าค่าคาดหวังของตัวดำเนินการแฮร์มิเชียนใดๆ จะต้องเป็นจำนวนจริงเสมอ ($\langle A \rangle \in \mathbb{R}$)
    \item กำหนดฟังก์ชันคลื่น $\psi(x) = C x e^{-\alpha x}$ สำหรับ $x \ge 0$ และเป็นศูนย์สำหรับ $x < 0$ จงหาค่า $C$ และคำนวณ $\langle x \rangle, \langle p_x \rangle$
    \item จงแสดงว่าตัวดำเนินการพาริตี $\hat{\Pi}$ ซึ่งนิยามโดย $\hat{\Pi}\psi(x) = \psi(-x)$ เป็นตัวดำเนินการแฮร์มิเชียนและยูนิแทรี พร้อมหาค่าเจาะจงที่เป็นไปได้
    \item จงพิสูจน์ว่าหาก $[\hat{H}, \hat{A}] = 0$ และ $\hat{A}$ ไม่ขึ้นกับเวลาอย่างชัดแจ้ง ค่าคาดหวัง $\langle A \rangle$ จะเป็นค่าคงที่ของการเคลื่อนที่ (Constant of Motion)
    \item ในกล่อง 1 มิติกว้าง $L$ จงคำนวณองค์ประกอบเมทริกซ์ของตำแหน่ง $x_{mn} = \braket{m}{\hat{x}}{n} = \int_0^L \psi_m^*(x) x \psi_n(x) dx$
    \item จงคำนวณหาการกระจายฟลักซ์ความน่าจะเป็นของคลื่นระนาบที่ตกกระทบขอบเขตศักย์ขั้นบันได
    \item จงแสดงว่าผลรวมของตัวดำเนินการฉายภาพ $\hat{P}_1 + \hat{P}_2$ จะเป็นตัวดำเนินการฉายภาพ ก็ต่อเมื่อ $\hat{P}_1 \hat{P}_2 = 0$
    \item จงเขียนสมการเศษตกค้างของ PINNs สำหรับสมการชเรอดิงเงอร์ที่มีสนามศักย์แบบฮาร์มอนิก $V(x) = \frac{1}{2}x^2$
    \item จงพิสูจน์ทฤษฎีบทไวรัล (Virial Theorem) $\langle T \rangle = \frac{1}{2}\langle x \frac{dV}{dx} \rangle$ โดยการพิจารณาอัตราการเปลี่ยนแปลงตามเวลาของตัวดำเนินการ $\hat{x}\hat{p}$
\end{enumerate}

% ==========================================================
% โครงงานจำลองเชิงคำนวณด้วย Python
% ==========================================================
\section{โครงงานจำลองเชิงคำนวณด้วยภาษา Python}

โปรแกรมภาษา Python ต่อไปนี้สาธิตการแก้สมการชเรอดิงเงอร์ไม่ขึ้นกับเวลา (TISE) ใน 1 มิติ โดยใช้ระเบียบวิธีผลต่างสืบเนื่องในการสร้างเมทริกซ์ฮามิลโทเนียนและหาค่าเจาะจง (Finite Difference Matrix Diagonalization)

\begin{lstlisting}[language=Python, caption={โปรแกรม Python หาค่าเจาะจงพลังงานและฟังก์ชันคลื่นด้วยระเบียบวิธีเมทริกซ์ผลต่างสืบเนื่อง}]
import numpy as np
import matplotlib.pyplot as plt
from scipy.linalg import eigh

# กำหนดพารามิเตอร์ทางฟิสิกส์ (Atomic units: hbar = 1, m = 1)
hbar = 1.0
m = 1.0

# กำหนดกริดเชิงพื้นที่
N = 1000               # จำนวนจุดกริด
L = 10.0               # ความกว้างของอาณาเขต [-L/2, L/2]
x = np.linspace(-L/2, L/2, N)
dx = x[1] - x[0]

# กำหนดสนามพลังงานศักย์ (ตัวอย่าง: ศักย์ฮาร์มอนิก V(x) = 0.5 * k * x^2)
k_spring = 2.0
V = 0.5 * k_spring * x**2

# สร้างเมทริกซ์อนุพันธ์อันดับสอง (Kinetic Energy Matrix) ด้วย Finite Difference 3 จุด
# d^2 psi / dx^2 approx (psi_{i+1} - 2*psi_i + psi_{i-1}) / dx^2
main_diag = (hbar**2 / (m * dx**2)) + V
off_diag = (-hbar**2 / (2.0 * m * dx**2)) * np.ones(N - 1)

# ประกอบเมทริกซ์ฮามิลโทเนียนสมมาตร (Tridiagonal Hamiltonian)
H = np.diag(main_diag) + np.diag(off_diag, -1) + np.diag(off_diag, 1)

# คำนวณค่าเจาะจงพลังงาน (Eigenvalues) และสถานะเจาะจง (Eigenvectors)
eigenvalues, eigenvectors = eigh(H)

# พล็อตระดับพลังงานและฟังก์ชันคลื่น 4 สถานะแรก
plt.figure(figsize=(10, 7))
plt.plot(x, V, 'k--', label='Potential $V(x)$', lw=1.5)

colors = ['#00F0FF', '#7B2CBF', '#FF007F', '#00FF66']
num_states = 4

for n in range(num_states):
    E_n = eigenvalues[n]
    psi_n = eigenvectors[:, n]
    
    # นอร์มอลไลซ์ฟังก์ชันคลื่น integral |psi|^2 dx = 1
    norm = np.sqrt(np.sum(psi_n**2) * dx)
    psi_n = psi_n / norm
    
    # วาดฟังก์ชันคลื่นโดยเลื่อนระดับขึ้นไปตามค่าพลังงาน E_n
    scale = 1.5
    plt.plot(x, E_n + scale * psi_n, color=colors[n], lw=2, 
             label=f'n={n}: $E_{n}={E_n:.3f}$')
    plt.axhline(E_n, color=colors[n], linestyle=':', alpha=0.5)

plt.xlim(-4.0, 4.0)
plt.ylim(0, eigenvalues[num_states-1] + 2.0)
plt.title('Quantum Harmonic Oscillator: Eigenenergies and Wavefunctions (Finite Difference)', 
          fontsize=13, fontweight='bold')
plt.xlabel('Position $x$', fontsize=12)
plt.ylabel('Energy / Shifted Wavefunction', fontsize=12)
plt.grid(True, linestyle=':', alpha=0.6)
plt.legend(loc='upper right', fontsize=10)
plt.tight_layout()
plt.savefig('schrodinger_finite_difference.png', dpi=300)
plt.show()
\end{lstlisting}
